\chapter{Inherited Parents and Recursive Formation}
\label{ch:inheritance}

A collection becomes a parent when its retained state supports inherited
interaction and evolution. A probability law obtained by compression is a first
part of this construction. It does not determine how that law will change, what
the parent can communicate, or which information its children must retain. We
give exact conditions for these further steps, two sufficient mechanisms
that make them hold, and a finite drive--response example in which exact
lumping works through two levels. The example supplies its reductions;
formation of the reductions remains a separate question.

Most exactness criteria below are instances of classical mathematics, and we
attribute them where they are stated: factor maps of dynamical systems, exact
lumping of differential equations, strong lumpability and Markov functions of
Markov chains, Riemannian submersions, equivariant reduction, and conditional
independence. What this chapter adds is their assembly into one notion of an
inherited parent with declared protocol and preparation classes, the
counterexamples showing which natural retentions fail, and the restrictions
that these criteria place on the formation constructions of later chapters.

\section{Two objects that a parent must retain}

\definitionheading{Sample compression and configuration retention}{inh:two-retentions}
Let $x\in\mathsf S$ be a complete agent configuration and let $Q_x$ be its
declared law on a standard Borel space $\mathsf Y$ of represented contents. A
normalized sample channel $C:\mathsf Y\rightsquigarrow\mathsf Z$ gives the
retained law
\begin{equation}
 Q_x^c(A)=\int C(y,A)Q_x(dy).
 \label{inh:sample-parent}
\end{equation}
A configuration retention is a measurable map
$R:\mathsf S\to\mathsf S^c$. Its value includes the parent laws and the
additional data needed for its declared interactions and update.
\Cref{inh:canonical-parent} identifies the smallest such data for a declared
interface.

The distinction applies to both sectors. A parent carrying $q_I$ and $s_I$
also needs the evaluator that makes $s_I$ predictive. Boundary channels,
source rows, source marks, and memory can be further components of $R(x)$.
Equal marginal laws do not imply equality of these components. When a sample
channel reads random links or source variables, those variables belong to the
joint input law of the channel. Substitution of a mean link or a selected
source is a different retention.

A law of represented contents and a probability law over complete
configurations are different objects. Coupled deterministic evolution of
law-valued agents requires no probability distribution over $x$. If such a
distribution is used below, its preparation or transition mechanism is stated.
All constructions admit general probability laws. A Gaussian parent is one
possible family, subject to its own closure conditions.

Throughout the chapter, $\circ$ composes maps and juxtaposition composes
kernels in the order in which they act: for kernels
$K:\mathsf X\rightsquigarrow\mathsf O$ and
$D:\mathsf O\rightsquigarrow\mathsf O'$,
$(KD)(x,A)=\int D(o,A)K(x,do)$. A measurable map used in a kernel product
stands for its deterministic kernel. With this convention every square below
is a commutative diagram of Markov kernels between standard Borel spaces, and
the composition results are the pasting of such diagrams in the Markov
category of these spaces \citep{Fritz2020}.

\section{Exact deterministic inheritance}

\lemmaheading{Factorization through a Borel surjection}{inh:borel-factorization}
Let $\mathsf S,\mathsf S^c,\mathsf Y$ be standard Borel spaces, let
$R:\mathsf S\to\mathsf S^c$ be a Borel surjection, and let
$h:\mathsf S\to\mathsf Y$ be Borel and constant on every fiber of $R$. There
is a unique Borel map $h^c:\mathsf S^c\to\mathsf Y$ with $h=h^c\circ R$.
\begin{proof}
Constancy on fibers and surjectivity define $h^c$ uniquely. Its graph is the
image of $\mathsf S$ under the Borel map $x\mapsto(R(x),h(x))$ into the
standard Borel space $\mathsf S^c\times\mathsf Y$, hence an analytic set. A
function between standard Borel spaces with analytic graph is Borel
\citep[Thm.~14.12]{Kechris1995}.
\end{proof}

The lemma is used below for maps into spaces of probability laws, which are
standard Borel when they carry the $\sigma$-algebra generated by the
evaluations $p\mapsto p(A)$ \citep{Kechris1995}.

\propositionheading{An update descends precisely when it respects retained states}{inh:deterministic-descent}
Let $\mathsf S,\mathsf S^c$ be standard Borel spaces, let
$R:\mathsf S\to\mathsf S^c$ be a Borel surjection, and let
$F:\mathsf S\to\mathsf S$ be Borel. There is a Borel parent update $F^c$
satisfying
\begin{equation}
 R\circ F=F^c\circ R,
 \label{inh:update-square}
\end{equation}
if and only if
\begin{equation}
 R(x)=R(x')\quad\Longrightarrow\quad R(F(x))=R(F(x')).
 \label{inh:fiber-test}
\end{equation}
The parent update is unique. The same conclusion holds on arbitrary
measurable spaces when $R$ has a measurable right inverse.

\begin{proof}
\Cref{inh:update-square} implies the implication in
\cref{inh:fiber-test}. Conversely, the implication says that $h=R\circ F$
is constant on the fibers of $R$, and \cref{inh:borel-factorization} gives a
unique Borel $F^c$ with $R\circ F=F^c\circ R$. With a measurable right
inverse $\iota$ on arbitrary measurable spaces, the formula
$F^c=R\circ F\circ\iota$ gives measurability directly, and surjectivity
gives uniqueness.
\end{proof}

This is the factor-map (semiconjugacy) criterion of dynamical systems. The
analytic-graph argument matters because a Borel surjection need not have a
Borel right inverse; the uniformization theorem supplies only universally
measurable ones \citep[Sec.~18.A]{Kechris1995}. Repeated substitution gives
$R\circ F^n=(F^c)^n\circ R$ for every nonnegative integer $n$. This is an
exact statement about the chosen update, including its order and retained
parameters. Replacing that update by descent of a newly defined parent
objective requires another argument.

For a differentiable version, suppose $R$ is continuously differentiable and
the fine and parent vector fields $B,B^c$ have unique local solutions. Their
flows satisfy $R\circ\Phi_t=\Phi_t^c\circ R$ on their common existence
interval if
\begin{equation}
 DR_x B(x)=B^c(R(x)).
 \label{inh:vector-square}
\end{equation}
Indeed, the chain rule makes $R(\Phi_t(x))$ a parent solution with initial
value $R(x)$, and uniqueness identifies it. Conversely, differentiation of
the flow identity at zero gives \cref{inh:vector-square}. For a moving map
$R_t$, the left side becomes $\partial_tR_t(x)+DR_t(x)B_t(x)$.
Membership changes therefore need an explicit handoff map; a differential
identity on one fixed partition does not specify them. Condition
\cref{inh:vector-square} is the exact-lumping condition for differential
equations, introduced for linear kinetics by \citet{WeiKuo1969}, for
nonlinear kinetics by \citet{LiRabitz1989}, and for flows on manifolds by
\citet{HorstmeyerAtay2016}; flow intertwining has also been proposed as a
criterion of macroscopic emergence \citep{BarnettSeth2023}. The
differentiable statement needs $R$ to be continuously differentiable between
manifolds on which both flows are defined. A retention into spaces of laws
that is differentiable only weakly, with a parent law of motion given by a
weak continuity equation, is not covered, and its intertwining must be
proved directly.

\propositionheading{Projectable vector fields}{inh:projectable}
Let $R:\mathsf S\to\mathsf S^c$ be a $C^{k+1}$ surjective submersion between
finite-dimensional manifolds, $k\geq1$, whose fibers are connected, and let
$B$ be a $C^k$ vector field on $\mathsf S$. A $C^k$ parent field $B^c$
satisfying \cref{inh:vector-square} exists if and only if, in every
submersion chart $R(a,b)=a$ with $B=(B_A,B_V)$,
\begin{equation}
 \partial_bB_A(a,b)=0.
 \label{inh:projectable-criterion}
\end{equation}
For a linear surjection $R=M$ on $\R^n$ this is the condition
$M(DB(x)v)=0$ for all $x$ and all $v\in\ker M$.

\begin{proof}
In a submersion chart, $DR_xB(x)=B_A(a,b)$. Condition
\cref{inh:projectable-criterion} says that this map is locally constant
along each fiber; a locally constant function on a connected fiber is
constant, so $x\mapsto DR_xB(x)$ is constant on fibers, and conversely
constancy on fibers gives \cref{inh:projectable-criterion}. A submersion
has smooth local sections $\sigma$ through every point, and
$B^c=(DR\cdot B)\circ\sigma$ is then well defined and $C^k$. Since $k\geq1$,
both fields have unique local solutions. In the linear case the fibers are
affine, hence connected, and the derivative of $MB$ along $v\in\ker M$ is
$M(DB(x)v)$.
\end{proof}

\remarkheading{Why inheritance is tested on a set of configurations}{inh:trajectory-remark}
Dynamical inheritance checked only along one realized trajectory carries no
information. If $x(t)$ is a smooth solution on $[0,T]$ and $R$ is smooth with
$z=R\circ x$ an embedded arc, the tangent field of the arc extends to a smooth
field $B^c$ on a neighborhood of the arc, and $R(\Phi_t x_0)=\Phi^c_t(Rx_0)$
holds with unique parent solutions. The same holds on a prepared family on
which $R$ is injective. The fiber test \cref{inh:fiber-test} constrains a
retention only over a set containing distinct configurations with the same
retained value. For this reason the inheritance identities are imposed below
on a declared invariant set $\Omega_I$ of post-handoff states
(\cref{inh:inherited-parent}), and the formation standard requires an
admitted class that is open in the chart topology of the dynamics or
relatively open in a declared invariant family that is not a finite union
of orbits.

\propositionheading{Gradient descent through a horizontally conformal submersion}{inh:gradient-inheritance}
Let $R:(\mathsf S,g)\to(\mathsf S^c,g_c)$ be a smooth surjective
submersion between finite-dimensional Riemannian manifolds, with horizontal
spaces $H_x=(\ker DR_x)^\perp$. Assume there is a smooth positive function
$d$ such that
\begin{equation}
 g_c(DR_xu,DR_xv)=d(x)^2g(u,v),
 \qquad u,v\in H_x.
 \label{inh:horizontal-conformality}
\end{equation}
Let $E,E_c$ be smooth potentials and $a$ a positive function. Then
\begin{equation}
 DR_x[-\operatorname{grad}_gE(x)]
 =a(x)[-\operatorname{grad}_{g_c}E_c](R(x))
 \label{inh:gradient-factor}
\end{equation}
holds at $x$ if and only if
\begin{equation}
 d(x)^2dE_x(v)=a(x)(dE_c)_{R(x)}(DR_xv)
 \qquad\text{for every }v\in H_x.
 \label{inh:horizontal-criterion}
\end{equation}
In particular, if $E=\chi\circ E_c\circ R$ with $\chi'>0$, then
\cref{inh:gradient-factor} holds with
$a(x)=\chi'(E_c(R(x)))d(x)^2$, and a Riemannian submersion with
$E=E_c\circ R$ gives same-time inheritance.
\begin{proof}
Let $w=\operatorname{grad}_gE(x)$ and let $w^H$ be its horizontal part.
Since $DR_xw=DR_xw^H$, for any coarse tangent $z$ with horizontal lift $v$,
\begin{equation}
 g_c(DR_xw,z)
 =d(x)^2g(w^H,v)
 =d(x)^2g(w,v)
 =d(x)^2dE_x(v).
\end{equation}
Nondegeneracy of $g_c$ shows that $DR_xw=a(x)\operatorname{grad}_{g_c}E_c(R(x))$
exactly when $d(x)^2dE_x(v)=a(x)(dE_c)_{R(x)}(z)$ for every $z$, which is
\cref{inh:horizontal-criterion} because horizontal lifting is a bijection
$T_{R(x)}\mathsf S^c\to H_x$. If $E=\chi\circ E_c\circ R$, then
$dE_x(v)=\chi'(E_c(R(x)))(dE_c)_{R(x)}(DR_xv)$ for every $v$. When $a=1$,
\cref{inh:vector-square} applies.
\end{proof}

This is the gradient pushforward for Riemannian submersions
\citep{ONeill1966} and horizontally conformal submersions
\citep{BairdWood2003}. The potential hypothesis $E=\chi\circ E_c\circ R$ is
much stronger than the criterion. It forces $dE$ to vanish on $\ker DR$, so
the fine gradient is horizontal and every fine trajectory is the horizontal
lift of its projection: for a product metric and $R(a,b)=a$, the discarded
coordinate does not move. Criterion \cref{inh:horizontal-criterion}
constrains only horizontal derivatives and leaves the discarded dynamics
free. On Euclidean $\R^2$ with $R(x_1,x_2)=(x_1+x_2)/\sqrt2$, the potential
\begin{equation}
 E=\tfrac12x_1^2+\tfrac12x_2^2+W(x_1-x_2),
 \qquad E_c(z)=\tfrac12z^2,
\end{equation}
with any smooth $W$, satisfies \cref{inh:horizontal-criterion} with $a=1$,
so the center of mass descends at the same time. With $s=x_1-x_2$ one has
$E=\tfrac12z^2+\tfrac14s^2+W(s)$, so $E\neq\chi\circ E_c\circ R$ unless
$W(s)=-\tfrac14s^2+\mathrm{const}$, and the separation obeys
$\dot s=-s-2W'(s)$. The metric hypothesis is necessary in a
uniform sense: if \cref{inh:gradient-factor} holds with $a\equiv1$ for every
potential of the form $E_c\circ R$, then
$DR_x\circ g^{-1}\circ DR_x^{*}=g_c^{-1}$ at $R(x)$, because
$DR(\operatorname{grad}_g(E_c\circ R))=(DR\circ g^{-1}\circ DR^{*})(dE_c)$ must equal
$g_c^{-1}dE_c$ for every covector, and this identity says that $R$ is a
Riemannian submersion. For a single potential the metric hypothesis is not
necessary.

For $a>0$, a fine trajectory projects onto the same oriented coarse
gradient orbit after the time change $d\tau/dt=a(x(t))$. An autonomous
field in the original time additionally requires the product
$a(x)\operatorname{grad}_{g_c}E_c(R(x))$ to be constant on retention
fibers. Factorization $a=\bar a\circ R$ suffices; dependence on discarded
state need not satisfy this requirement. The proposition is one sufficient
route to inheritance. Fisher contraction and preservation of the
collective connection in \cref{ch:geometry} are different conditions.

Fiber minimization alone does not supply the criterion.
For $R(u,v)=v$ on Euclidean $\R^2$, let
\begin{equation}
 E(u,v)=\tfrac12(u-v)^2+\tfrac12v^2,
 \qquad E_c(v)=\inf_u E(u,v)=\tfrac12v^2.
\end{equation}
Equality $E=E_c\circ R$ holds only on $u=v$. The horizontal direction is
$\partial_v$, and $dE(\partial_v)-dE_c(DR(\partial_v))=v-u$, so the
same-time criterion fails off the diagonal; the pushed field $u-2v$ is not a
function of $v$. Since $dE(\partial_u)=u-v$ does not vanish on any open set,
there is also no identity $E=\chi(E_c\circ R)$ on an open set.

\propositionheading{Equivariant descent}{inh:equivariant-descent}
Let a Lie group $G$ act smoothly, freely, and properly on a
finite-dimensional manifold $\mathsf S$, and let
$\pi:\mathsf S\to\mathsf S/G$ be the quotient map.
\begin{enumerate}
\item[(i)] If a smooth vector field $B$ is $G$-equivariant, there is a unique
smooth field $B^c$ on $\mathsf S/G$ with $D\pi\cdot B=B^c\circ\pi$, and the flows
intertwine. Every $G$-invariant measurable response kernel or channel on
$\mathsf S$ factors through $\pi$.
\item[(ii)] If $B=-\operatorname{grad}_gE$ with $g$ and $E$ both
$G$-invariant, then $\mathsf S/G$ carries a unique metric $g_c$ making $\pi$
a Riemannian submersion, $E=E_c\circ\pi$ for a smooth $E_c$, and
$B^c=-\operatorname{grad}_{g_c}E_c$.
\end{enumerate}
\begin{proof}
By the quotient manifold theorem \citep[Thm.~21.10]{Lee2013}, $\mathsf S/G$
is a smooth manifold and $\pi$ a surjective submersion whose fibers are the
orbits. Since $\pi\circ\phi_h=\pi$ for the action $\phi_h$,
$D\pi_{hx}D\phi_h=D\pi_x$, and equivariance gives
$D\pi_{hx}B(hx)=D\pi_{hx}D\phi_hB(x)=D\pi_xB(x)$. Thus $D\pi\cdot B$ is constant
on fibers, and $B^c=(D\pi\cdot B)\circ\sigma$ for local smooth sections $\sigma$ is
well defined and smooth. Invariant kernels are constant on orbits and
descend by \cref{inh:deterministic-descent} applied to the kernel-valued map.
For (ii), invariance of $E$ makes $dE$ vanish on orbit tangents, which for a
free action span $\ker D\pi$. Isometries carry horizontal spaces to
horizontal spaces, so $g_c(D\pi u,D\pi v)=g(u,v)$ for horizontal $u,v$ is well
defined \citep{ONeill1966}. \Cref{inh:gradient-inheritance} with $d=1$ and
$\chi=\id$ gives the gradient identity.
\end{proof}

Equivariant descent holds on the whole configuration space, so the fiber
test needs no restriction to an admitted class. It forgets only orbit data.
Without freeness the quotient is stratified by orbit type and the statement
holds on each stratum. On infinite-dimensional law charts, (ii) requires in
addition a closed complement to the orbit tangent in the Fisher metric and a
quotient Banach manifold, which must be verified for the chart in use.

\propositionheading{Coordinate projections are exact only for isolated blocks}{inh:coordinate-projection}
Let $\mathsf S=\mathsf S_A\times\mathsf S_B$ be a product of connected
manifolds with metric $g=g_A(a)\oplus g_B(b)$, and let $E\in C^2$. The
projection $R(a,b)=a$ intertwines $-\operatorname{grad}_gE$ with an
autonomous field if and only if $\partial_a\partial_bE\equiv0$, that is,
$E(a,b)=E_A(a)+E_B(b)$. The discarded block is then autonomous as well.
\begin{proof}
The retained velocity is $\dot a=-g_A(a)^{-1}\partial_aE(a,b)$. Since the
fibers $\{a\}\times\mathsf S_B$ are connected and $g_A(a)$ is invertible,
it is independent of $b$ exactly when $\partial_b\partial_aE=0$, that is,
when $\partial_aE(a,b)=\phi(a)$. Integrating along paths in the connected
factors gives $E=E_A(a)+E_B(b)$, and then
$\dot b=-g_B(b)^{-1}\partial_bE_B(b)$.
\end{proof}

The Fisher flow of \cref{ch:dynamics} with product recognition has
agent-local metric blocks and one potential, and a source appears in the
cost of every receiver that attends to it. Discarding whole agents or whole
sectors is therefore exact in that dynamics only for energetically isolated
blocks, with no peer cost in either direction between the retained and the
discarded agents. Exact compression of reciprocally coupled agents can come
from symmetry (\cref{inh:equivariant-descent}) or from horizontal energy
structure (\cref{inh:gradient-inheritance}), as in the center-of-mass example.
These are sufficient mechanisms, not a classification; the general criterion
is projectability (\cref{inh:projectable}).

\propositionheading{Exact inheritance of a gradient flow creates no oscillatory mode at images of critical points}{inh:no-underdamped}
Let $B=-\operatorname{grad}_gE$ with $E\in C^2$, and let $x^*$ be a critical
point of $E$. Then $L=DB(x^*)=-g(x^*)^{-1}\operatorname{Hess}E(x^*)$ is
self-adjoint for $g(x^*)$ and has real spectrum. If a $C^2$ submersion $R$
intertwines $B$ with a $C^1$ field $B^c$ near $x^*$ in the sense of
\cref{inh:vector-square}, then $R(x^*)$ is an equilibrium of $B^c$ and
$\operatorname{spec}DB^c(R(x^*))\subseteq\operatorname{spec}L\subset\R$.
If $z^*$ is an equilibrium of $B^c$ and $E$ restricted to the fiber
$R^{-1}(z^*)$ attains its infimum, for example on a compact fiber, then the
fiber contains a critical point of $E$ and the conclusion holds at $z^*$.
\begin{proof}
For the last statement, $B^c(z^*)=0$ and \cref{inh:vector-square} make
$B(x)$ vertical at every $x\in R^{-1}(z^*)$, so $dE_x$ vanishes on the
horizontal space $(\ker DR_x)^\perp$. At a minimizer of $E$ on the fiber,
$dE_x$ also vanishes on the tangent space $\ker DR_x$ of the fiber, so the
minimizer is a critical point. For the main statement:
At a critical point the Hessian is a symmetric bilinear form, so $L$ is
$g(x^*)$-self-adjoint. Evaluating \cref{inh:vector-square} at $x^*$ gives
$B^c(R(x^*))=0$. Differentiating it at $x^*$, the term containing the second
derivative of $R$ is multiplied by $B(x^*)=0$, leaving
$DR\circ L=L^c\circ DR$ with $L^c=DB^c(R(x^*))$. Hence $\ker DR$ is $L$-invariant,
and since $DR$ is surjective, $L^c$ is conjugate through $DR$ to the map that
$L$ induces on $T_{x^*}\mathsf S/\ker DR$. The characteristic polynomial of
$L$ is the product of those of its restriction and of the induced quotient
map, so every eigenvalue of $L^c$ is an eigenvalue of $L$.
\end{proof}

An exactly inherited parent of a Fisher gradient flow therefore has no
underdamped linear mode at the image of a critical point. The critical-point
hypothesis cannot be dropped. Let one agent have law $\mathcal N(z,I_2)$ on
$\R^2$, with Euclidean Fisher metric, and another the Poisson law of mean
$e^{-y}$, with Fisher metric $e^{-y}dy^2$, and take
$E=\tfrac\mu2|z|^2+e^{-y}\geq0$ with $\mu>0$. Then $E$ has no critical
point, the flow is $\dot z=-\mu z$, $\dot y=1$, and the retention
$w=\operatorname{Rot}(\kappa y)z$, which reads the first agent in a frame
turned by the second, is exact with parent field
$\dot w=(-\mu I+\kappa J)w$, $J$ the rotation generator, whose eigenvalues
$-\mu\pm i\kappa$ are not real. Oscillatory collective modes therefore
require nonreciprocal fields, explicit driving, declared second-order
kinetics, or parent equilibria over fibers along which the fine state
drifts without reaching a critical point.

\propositionheading{Loss of dynamical roles under product recognition}{inh:role-loss}
Take $Q_x=\operatorname{Ber}(p_1)\otimes\operatorname{Ber}(p_2)$, with
$x=(p_1,p_2)\in[0,1]^2$, and the coupled deterministic law
\begin{equation}
 \dot p_1=-p_1+\lambda(p_2-p_1),\qquad
 \dot p_2=-2p_2+\lambda(p_1-p_2),\qquad \lambda\geq0.
 \label{inh:role-flow}
\end{equation}
Retain only the law of $K_1+K_2$. No autonomous update of that retained law
reproduces every trajectory of \cref{inh:role-flow}.

\begin{proof}
The linear system has a unique global solution. Its vector field does not
point outward on any face of the square: at $p_1=0$, $\dot p_1=\lambda p_2\geq0$,
at $p_1=1$, $\dot p_1\leq-1$, and similarly for $p_2$. The square is therefore
forward invariant; the field is tangent to a face when $\lambda=0$ and
vanishes at the corner $(0,0)$. The configurations $(0.8,0.2)$ and
$(0.2,0.8)$ both give probabilities $(0.16,0.68,0.16)$ for the values
$0,1,2$. However,
\begin{equation}
 \frac{d}{dt}\E(K_1+K_2)=-(p_1+2p_2),
\end{equation}
equals $-1.2$ and $-1.8$ respectively. Their retained laws therefore have
different derivatives, and no autonomous update of the retained law can
reproduce both trajectories.
\end{proof}

This obstruction occurs while recognition remains a product. The missing
datum is which probability belongs to which dynamical role. Correlated
recognition is neither a prerequisite for this counterexample nor its
automatic remedy. Coupled mean-field dynamics are admissible; the proposed
retention must still pass its own inheritance test.

\section{Exact stochastic inheritance}

\propositionheading{Finite strong lumpability}{inh:finite-lumpability}
Let $T$ be a transition matrix on a finite configuration space $\mathsf S$
and let $R:\mathsf S\to\mathsf S^c$ be surjective. A single parent
transition matrix reproduces the retained process for every initial law if
and only if, for every $z'\in\mathsf S^c$,
\begin{equation}
 \sum_{y:R(y)=z'}T(x,y)
 =\sum_{y:R(y)=z'}T(x',y)
 \quad\text{whenever }R(x)=R(x').
 \label{inh:lumpability}
\end{equation}
In that case the common sum defines $T^c(R(x),z')$ uniquely.

\begin{proof}
The common sums are nonnegative and sum to one, so they define a transition
matrix. Conditional on any retained history ending at $z$, the hidden
configuration lies in $R^{-1}(z)$. Every state in that fiber gives the same
next retained distribution $T^c(z,\cdot)$; averaging over its conditional
law leaves this distribution unchanged. Induction gives the parent path law
from the pushed initial law. Conversely, start the fine process at each
point mass $\delta_x$. A common parent kernel must give the same one-step
distribution whenever these initial states have the same image. This is
\cref{inh:lumpability}.
\end{proof}

This is the finite strong-lumpability criterion
\citep{KemenySnell1976}; its general state-space form is Dynkin's criterion
for a function of a Markov process \citep{Dynkin1965}. On standard Borel
spaces the same argument applies whenever $x\mapsto T(x,R^{-1}(\cdot))$ is
constant on fibers: the analytic-graph argument of
\cref{inh:deterministic-descent}, applied to this map into the standard Borel
space of probability laws on $\mathsf S^c$, makes the parent kernel
measurable. Closure for every initial law is stronger than closure for one
chosen preparation. The latter is weak lumpability \citep{KemenySnell1976},
and the question whether a function of a Markov chain is itself Markov goes
back to \citet{BurkeRosenblatt1958}; \citet{Buchholz1994} compares these
notions with exact lumpability.

An equilibrium conditional average of fine transitions defines a normalized
candidate parent kernel. It does not prove \cref{inh:lumpability}: the
hidden conditional law during a transient can differ from the equilibrium
conditional. Likewise, agreement of invariant distributions does not
establish agreement of histories.

\propositionheading{Intertwined handoffs}{inh:intertwining}
Let $T$ be a transition matrix on a finite space $\mathsf S$, let
$R:\mathsf S\to\mathsf S^c$ be surjective, and let $\Phi_R$ be its
deterministic kernel. Let $\Lambda$ be a stochastic matrix from
$\mathsf S^c$ to $\mathsf S$ with $\Lambda(z,R^{-1}\{z\})=1$ for every $z$,
and suppose
\begin{equation}
 \Lambda T=T^c\Lambda
 \label{inh:intertwining-identity}
\end{equation}
for some matrix $T^c$. Then $T^c=\Lambda T\Phi_R$ is a transition matrix.
For every law $\nu$ on $\mathsf S^c$, under the initial law $\nu\Lambda$ the
retained process $Z_n=R(X_n)$ is Markov with initial law $\nu$ and kernel
$T^c$, and on every retained history of positive probability the conditional
law of $X_n$ is $\Lambda(Z_n,\cdot)$. For a response kernel $K_a$ on
$\mathsf S$, the predictive response given the retained history is
$(\Lambda K_a)(Z_n,\cdot)$.

\begin{proof}
The support condition gives $\Lambda\Phi_R=I$, so multiplying
\cref{inh:intertwining-identity} on the right by $\Phi_R$ yields
$T^c=\Lambda T\Phi_R$, a product of stochastic matrices. Let $P_z$ be the
diagonal selector of $R^{-1}\{z\}$. The support condition also gives
$\Lambda(w,\cdot)P_z=\delta_{wz}\Lambda(z,\cdot)$. Let $f_n$ be the row
vector whose $x$ entry is the probability of the history
$(z_0,\ldots,z_n)$ jointly with $X_n=x$. Then
$f_0=\nu\Lambda P_{z_0}=\nu(z_0)\Lambda(z_0,\cdot)$. If
$f_n=c_n\Lambda(z_n,\cdot)$, then
\begin{equation}
 f_{n+1}=f_nTP_{z_{n+1}}
 =c_n(T^c\Lambda)(z_n,\cdot)P_{z_{n+1}}
 =c_nT^c(z_n,z_{n+1})\Lambda(z_{n+1},\cdot).
\end{equation}
By induction $c_n=\nu(z_0)\prod_{k<n}T^c(z_k,z_{k+1})$ is the probability
of the history and $\Lambda(z_n,\cdot)$ is the normalized conditional law.
The predictive response follows by integrating $K_a$ against it.
\end{proof}

This is the Markov-function theorem of \citet{RogersPitman1981}, who also
give the continuous-time and general state-space versions. It gives a
sufficient condition for a handoff ``from a declared conditional law'' to
yield a memoryless parent, namely the kernel $\Lambda$ and an initial law of
the form $\nu\Lambda$, and it identifies the filter as $\Lambda(Z_n,\cdot)$.
It is not necessary: a strongly lumpable chain gives a memoryless parent
from every initial law, and on $\{1,2,3\}$ with rows $(0.6,0,0.4)$,
$(0,0.6,0.4)$, $(0.5,0,0.5)$, $R(1)=R(2)=a$, $R(3)=b$, the chain is strongly
lumpable while among $\Lambda(a,\cdot)=(\lambda,1-\lambda,0)$,
$\Lambda(b,\cdot)=\delta_3$ the identity \cref{inh:intertwining-identity}
holds only at $\lambda=1$. Intertwining does not require strong
lumpability. On
$\mathsf S=\{1,2,3\}$ with $R(1)=R(2)=a$, $R(3)=b$, rows
$(0.3,0.5,0.2)$, $(0.4,0.2,0.4)$, $(0.25,0.25,0.5)$, and
$\Lambda(a,\cdot)=(\tfrac12,\tfrac12,0)$, $\Lambda(b,\cdot)=\delta_3$,
condition \cref{inh:lumpability} fails ($0.2\neq0.4$) while
\cref{inh:intertwining-identity} holds with
$T^c=\bigl(\begin{smallmatrix}7/10&3/10\\1/2&1/2\end{smallmatrix}\bigr)$.
From the initial law $\delta_1$, which is not of the form $\nu\Lambda$, the
retained path laws of length five differ from those of $T^c$ by $1/10$ in
total variation.

\propositionheading{The coarsest exact parent containing an interface}{inh:canonical-parent}
Let $F:\mathsf S\to\mathsf S$ and an interface $Q:\mathsf S\to\mathsf Y$ be
measurable. The interface collects what the parent must reproduce: its laws,
the responses of a declared protocol class, and the channels and evaluators
it reads. Define
\begin{equation}
 \pi(x)=(Q(F^nx))_{n\geq0}\in\mathsf Y^{\mathbb N}.
 \label{inh:canonical-retention}
\end{equation}
Then $\pi$ satisfies \cref{inh:fiber-test}, with parent update the left
shift on $\pi(\mathsf S)$, and $Q$ is its zeroth coordinate. If a retention
$R$ satisfies \cref{inh:fiber-test} and $Q=Q^c\circ R$ for some $Q^c$, then
$R(x)=R(x')$ implies $\pi(x)=\pi(x')$. Thus $\pi$ is, up to relabeling of its
values, the unique coarsest exact retention through which the interface
factors. For a forward-complete continuous flow on a Polish space and a
continuous $Q$ into a Polish space $\mathsf Y$, replace $F^n$ by $\Phi_t$,
$t\geq0$, and take values in the Polish space $C([0,\infty),\mathsf Y)$. If
$B$ and $Q$ are real analytic, $\pi(x)=\pi(x')$ exactly when
$L_B^kQ(x)=L_B^kQ(x')$ for every $k\geq0$. The image $\pi(\mathsf S)$ is in
general an analytic set, and $\pi$ is a configuration retention in the
sense of \cref{inh:two-retentions}, onto a standard Borel parent space, when
$\pi(\mathsf S)$ is Borel; this holds, for example, when $\mathsf S$ is
countable or $\pi$ is injective on a Borel transversal of its fibers.

\begin{proof}
The identity $\pi(F(x))=(Q(F^{n+1}x))_{n\geq0}$ is the shift of $\pi(x)$, so
equal values have equal successors. If $R$ satisfies \cref{inh:fiber-test},
induction gives $R(F^nx)=R(F^nx')$ whenever $R(x)=R(x')$, and applying $Q^c$
gives $\pi(x)=\pi(x')$. In the analytic case $t\mapsto Q(\Phi_tx)$ is real
analytic on its interval of existence and its derivatives at zero are the
Lie derivatives $L_B^kQ(x)$, so the identity theorem gives the stated
equivalence.
\end{proof}

The analytic case is the observability construction of
\citet{HermannKrener1977}. For finite Markov chains the corresponding object
is computed by partition refinement.

\propositionheading{The coarsest strongly lumpable refinement with responses}{inh:canonical-lumping}
Let $T$ be a transition matrix on a finite space $\mathsf S$, let
$\mathcal B_0$ be the partition of $\mathsf S$ into level sets of an
interface $Q$, and let finitely many protocols $a$ have finite retained
response alphabets $\mathsf O_a^c$ and joint kernels
$J_a(x,\{o\}\times\{y\})=\Pr(D_a(O_a)=o,X'=y\mid X=x)$. Define
$\mathcal B_{k+1}$ by splitting each block of $\mathcal B_k$ according to
the vector of masses $J_a(x,\{o\}\times B)$ over all protocols $a$, values
$o\in\mathsf O_a^c$, and blocks $B\in\mathcal B_k$.
The sequence stabilizes at a partition $\mathcal B_*$, which is the unique
coarsest partition refining $\mathcal B_0$ that satisfies
\cref{inh:lumpability} and on which every joint kernel of retained response
and next block factors through the block.
\begin{proof}
Each step refines the previous partition, so the sequence stabilizes after
finitely many steps. At the fixed point, states in one block have equal
joint masses into every pair of a response value and a block; summing over
$o$ gives \cref{inh:lumpability}. Let $\mathcal B$ be any partition with
these properties refining $\mathcal B_0$, and suppose it refines
$\mathcal B_k$. Each block of $\mathcal B_k$ is a union of blocks of
$\mathcal B$, so two states in one block of $\mathcal B$ have equal joint
masses into every pair $(o,B)$ with $B\in\mathcal B_k$ and lie in one block
of $\mathcal B_{k+1}$. By induction $\mathcal B$ refines $\mathcal B_*$.
\end{proof}

With no protocols this is probabilistic bisimulation \citep{LarsenSkou1991},
and with them it is bisimulation with outputs; an efficient algorithm is
given by \citet{DerisaviHermannsSanders2003}. The joint kernels are needed:
the response example after \cref{inh:boundary-square} has a strongly
lumpable coordinate partition whose joint response kernel does not factor,
and splitting by the marginal response laws alone would return it. The two
propositions make the phrase ``the additional data needed'' in
\cref{inh:two-retentions} a theorem for deterministic updates and for finite
stochastic updates with preparation class equal to all initial laws: given
the update and the interface, the canonical parent is determined. For a
restricted preparation class no coarsest exact parent is stated here. A
larger interface gives a finer canonical parent, so the protocol class
decides how much a parent can forget. Unrestricted identity-channel probes
make it injective on exposed child laws
(\cref{sec:public-gibbs-identifiability}, \cref{prop:boundary-observable}).
Every exact retention whose interface contains the declared protocols
refines the canonical parent, so noninjectivity of a handoff retention
implies noninjectivity of the canonical parent but not conversely. The
compression clause (F4) of the formation standard concerns the handoff
retention itself and the child-level probe class of
\cref{inh:protocol-classes}.

\section{Boundary responses and composition}

A parent must inherit what it does to the rest of the population. Fix a
declared boundary protocol $a$, including the preparation and receiving
conditions. Let $K_a:\mathsf S\rightsquigarrow\mathsf O_a$ be its
normalized response kernel and let
$D_a:\mathsf O_a\rightsquigarrow\mathsf O_a^c$ retain the response of
interest. An inherited kernel $K_a^c$ satisfies
\begin{equation}
 (K_aD_a)(x,A)=K_a^c(R(x),A).
 \label{inh:boundary-square}
\end{equation}
On standard Borel spaces, such a kernel exists exactly when
$x\mapsto(K_aD_a)(x,\cdot)$ is constant on each fiber of $R$; the
analytic-graph argument of \cref{inh:deterministic-descent} gives
measurability, and substitution proves necessity. When several outputs are
used jointly, this condition applies to their joint kernel. Preserving each
output marginal alone need not preserve their dependence.

The same applies to a response and the next retained state. For a
stochastic update, response inheritance concerns the joint kernel
$x\mapsto\operatorname{law}(D_a(O_a),R(X'))$ on
$\mathsf O_a^c\times\mathsf S^c$, which must factor through $R$. Separate
descent of the response kernel and of the retained transition does not imply
it. Let $x=(u,v)\in\{0,1\}^2$, $R(x)=u$, draw a fresh fair bit $\xi$, set
$u'=u\oplus\xi$, draw $v'$ freshly, and respond with $O=\xi\oplus v$. Both
the response and the next retained state are uniform from every
configuration, but at $(0,0)$ the joint law of $(O,u')$ is uniform on
$\{(0,0),(1,1)\}$ and at $(0,1)$ it is uniform on $\{(1,0),(0,1)\}$. If the
randomness of the response is conditionally independent of the transition
randomness given $x$, separate descent suffices. For adaptive protocols, in
which a receiver chooses its next action from past outputs, the joint
condition is required for each admitted action.

\propositionheading{Exact elimination through a conditional separator}{inh:separator-gluing}
Let $Z,U,V$ have a declared joint reference law
$\omega(dz,du,dv)=\omega_Z(dz)\omega_{UV}(du,dv\mid z)$ on standard Borel
spaces. Denote its conditional marginals by $\omega_U,\omega_V$. For
bounded nonnegative likelihood factors $L_U(u,z),L_V(v,z)$, jointly
measurable, define
\begin{equation}
 F_U(z)=\int L_U(u,z)\omega_U(du\mid z),
 \qquad F_V(z)=\int L_V(v,z)\omega_V(dv\mid z).
\end{equation}
The identity
\begin{equation}
 \int L_U(u,z)L_V(v,z)\omega_{UV}(du,dv\mid z)
 =F_U(z)F_V(z)
 \label{inh:separator-factorization}
\end{equation}
holds $\omega_Z$-almost everywhere for every such pair precisely when
$U$ and $V$ are conditionally independent given $Z$.
Under that condition, a nonnegative factor $L_Z(z)$ with finite positive
total evidence gives the unnormalized retained measure
\begin{equation}
 m_Z(dz)=L_Z(z)F_U(z)F_V(z)\omega_Z(dz).
 \label{inh:separator-measure}
\end{equation}
\begin{proof}
Conditional independence factors $\omega_{UV}$ as
$\omega_U\otimes\omega_V$ almost surely. Tonelli's theorem then gives
\cref{inh:separator-factorization,inh:separator-measure}. Conversely,
take indicator factors from countable generating families closed under
finite intersections for the two standard Borel spaces. Outside one
$\omega_Z$-null set, the conditional joint and product laws agree on every
generating rectangle. Uniqueness of probability measures extends this
equality to the product sigma algebra, proving conditional independence.
\end{proof}

The shared factor $L_Z$ appears once. If the hidden blocks are not
conditionally independent, their joint conditional factor must be
retained. This elimination statement concerns the declared reference and
likelihood; it does not replace the law-dependent interaction potential
by a fixed joint model. It is the separator factorization of
\citet{Dawid1979} that underlies junction-tree propagation
\citep{Lauritzen1996}.

The proposition presupposes the joint law. A gluing statement holds for
two blocks: laws of $(Z,U)$ and $(Z,V)$ with a common $Z$ marginal have
exactly one joint law with $U$ and $V$ conditionally independent given $Z$,
namely the conditional product. Elimination of $k\geq3$ blocks needs mutual
conditional independence, not pairwise independence: with $U,V$ independent
fair bits, $W=U\oplus V$, and $Z$ constant, every pair is independent, but
$\E[1_{U=1}1_{V=1}1_{W=1}]=0\neq1/8=F_UF_VF_W$. Gluing over a family of
overlapping separators requires a junction tree. Edge laws on a four-cycle
forcing $A=B$, $B=C$, $C=D$, and $D\neq A$ with uniform single-site marginals
are pairwise consistent but have no joint law; existence of a joint law for
every consistent family characterizes the decomposable (running-intersection)
hypergraphs \citep{Vorobev1962}.

For a child channel $\mathcal K_{ij}:\mathsf Y_j\rightsquigarrow\mathsf Y_i$
and retentions $C_j:\mathsf Y_j\rightsquigarrow\mathsf Z_J$ and
$C_i:\mathsf Y_i\rightsquigarrow\mathsf Z_I$, the corresponding square is
\begin{equation}
 \mathcal K_{ij}C_i=C_j\mathcal K_{IJ}^c,
 \qquad \mathcal K_{IJ}^c:\mathsf Z_J\rightsquigarrow\mathsf Z_I.
 \label{inh:channel-square}
\end{equation}
The right side means first retain the sender and then apply the parent
channel. A parent channel exists exactly when $\mathcal K_{ij}C_i$ is a
garbling of $C_j$, that is, when it lies in the set $C_j\mathcal M$ with
$\mathcal M$ ranging over Markov kernels $\mathsf Z_J\rightsquigarrow\mathsf Z_I$;
for finite alphabets this is a linear feasibility problem. For a
deterministic retention $C_j$ it reduces to constancy of
$\mathcal K_{ij}C_i$ on the fibers of $C_j$: the retained sender must
determine every distinction visible to the retained receiver. For a
stochastic $C_j$ the condition is stronger than any linear solvability. If
$C_j$ is the binary symmetric channel with flip probability $1/4$ and
$\mathcal K_{ij}=C_i=\id$, the only linear solution is
$C_j^{-1}=\bigl(\begin{smallmatrix}3/2&-1/2\\-1/2&3/2\end{smallmatrix}\bigr)$,
which is not a Markov kernel, so no parent channel exists. For a
law-dependent channel $\mathcal K_{ij}[x]$, the square must hold at each
configuration and the map $x\mapsto\mathcal K_{ij}[x]C_i$ must also factor
through the configuration retention. In a deterministic linear realization
the square reads $C_i\circ M_{ij}=M_{IJ}^c\circ C_j$; the kernel-inclusion
criterion and its proof are given in the supplement.

\propositionheading{Inheritance composes}{inh:composition}
Suppose $R_{10}:\mathsf S_0\to\mathsf S_1$ and
$R_{21}:\mathsf S_1\to\mathsf S_2$ intertwine deterministic updates at
successive levels. Then $R_{20}=R_{21}\circ R_{10}$ intertwines the first
and third updates. The analogous assertion holds for finite lumped Markov
transitions. Let $K_a^\ell:\mathsf S_\ell\rightsquigarrow\mathsf O_\ell$
be boundary kernels, with response retentions
$D_{10}:\mathsf O_0\rightsquigarrow\mathsf O_1$ and
$D_{21}:\mathsf O_1\rightsquigarrow\mathsf O_2$. If
\begin{equation}
 K_a^0D_{10}=R_{10}K_a^1,\qquad
 K_a^1D_{21}=R_{21}K_a^2,
 \label{inh:two-boundaries}
\end{equation}
then $K_a^0D_{10}D_{21}=R_{10}R_{21}K_a^2$. The same pasting composes the
joint kernels $J^\ell_a$ of retained response and next retained state of
clause (B) below: $J^0_a(D_{10}\otimes R_{10})=R_{10}J^1_a$ and
$J^1_a(D_{21}\otimes R_{21})=R_{21}J^2_a$ give
$J^0_a(D_{10}D_{21}\otimes R_{10}R_{21})=R_{10}R_{21}J^2_a$. Intertwinings
compose: if $\Lambda_{10}T^0=T^1\Lambda_{10}$ and
$\Lambda_{21}T^1=T^2\Lambda_{21}$ with supports on the respective fibers,
then $(\Lambda_{21}\Lambda_{10})T^0=T^2(\Lambda_{21}\Lambda_{10})$, and
$\Lambda_{21}\Lambda_{10}(z,\cdot)$ is supported on the fiber of
$R_{21}\circ R_{10}$ over $z$. Canonical parents compose along nested
interfaces: if $Q_2=Q_{21}\circ Q_1$, the canonical parent of
\cref{inh:canonical-parent} for $Q_2$ is the canonical parent, for the
interface $Q_{21}$ applied to the zeroth coordinate, of the shift on
$\pi_{Q_1}(\mathsf S)$, composed with $\pi_{Q_1}$.
\begin{proof}
For deterministic updates,
$R_{21}\circ R_{10}\circ F^0=R_{21}\circ F^1\circ R_{10}
=F^2\circ R_{21}\circ R_{10}$. For Markov transitions and their
deterministic retention kernels, in kernel order,
$T^0R_{10}R_{21}=R_{10}T^1R_{21}=R_{10}R_{21}T^2$, and the kernel
$R_{10}R_{21}$ is the deterministic kernel of $R_{21}\circ R_{10}$.
For boundary kernels and joint kernels, associate the successive integrals
and substitute \cref{inh:two-boundaries} or the joint identities. All
integrands for events are nonnegative, so interchanging the integration
order is valid. For intertwinings,
$\Lambda_{21}\Lambda_{10}T^0=\Lambda_{21}T^1\Lambda_{10}
=T^2\Lambda_{21}\Lambda_{10}$, and the support statement follows because
$\Lambda_{21}(z,\cdot)$ charges only $R_{21}^{-1}\{z\}$ and each
$\Lambda_{10}(z',\cdot)$ charges only $R_{10}^{-1}\{z'\}$. For canonical
parents, both maps identify $x,x'$ exactly when $Q_2(F^nx)=Q_2(F^nx')$ for
every $n\geq0$.
\end{proof}

When changes of gauge presentation act on the retained spaces, the
retentions and boundary kernels must transform compatibly with those
actions. Re-expressing a commuting square then gives the same square in
the new coordinates. Neither a parent frame chosen for this purpose nor
the square itself determines a base connection. If the parent is also
required to retain the collective Fisher connection, the additional
condition in \cref{geo:retained-connection} must be checked.

\definitionheading{Declared protocol and probe classes}{inh:protocol-classes}
A construction declares, as part of its constitutive data, two rules defined
for every block $I$ of every partition of its agents at every level,
singletons included.
\begin{enumerate}
\item[(A)] The parent protocol rule assigns to $I$ a class $A_I$ of boundary
protocols. Each $a\in A_I$ consists of its preparation and receiving
conditions, its admitted receiver actions, its normalized response kernel
$K_a$, and its response retention $D_a$, and it addresses $I$ only through
channels and records that read the block as a whole; it does not name a
member.
\item[(P)] The child-level probe rule assigns to $I$ a class $B_I$ of
probes that some agent, channel, record kernel, or declared external
querier of the construction can apply to individual members of $I$, or to
sets of them, before a handoff of $I$. It contains every protocol of $A_J$
for every block $J$ merged into $I$ at that handoff, applied before the
handoff. A probe $b$ has a response law $\operatorname{law}_x(b)$ at each
configuration $x$.
\end{enumerate}
Both rules are equivariant under relabeling of the agents and are functions
of the index set of the block and of the constitutive data alone. They do
not depend on the retention, on which blocks are selected, or on the
realized trajectory. This is the content of the requirement that the
classes be fixed before the retention: a class defined from a retention,
such as the protocols whose responses factor through $R_I$ or the responses
computed from a retained law, is not a declared class. The probes of $B_I$
are elements of the observable algebra of \cref{ch:foundations} at the
child level, so a distinction that no $b\in B_I$ separates is gauge relative
to $B_I$ in the sense of that chapter. A protocol whose response retention
is constant imposes nothing in (B) below; a rule consisting of such
protocols is a declared rule and is reported as such.

\definitionheading{Inherited parent}{inh:inherited-parent}
An inherited parent of a block $I$ selected at a handoff consists of the
parent protocol class $A_I$ of \cref{inh:protocol-classes}; for stochastic
dynamics, a preparation class $\mathcal P_I$ of initial laws that contains
every post-handoff law the handoff produces from the admitted class, together
with any kernel $\Lambda$ used to define it; a Borel invariant set $\Omega_I$
that contains the reachable post-handoff set $\Omega_I^{\mathrm{reach}}$, the
forward orbits of the post-handoff configurations arising from the admitted
class with the block $I$; and a Borel retention $R_I:\Omega_I\to\mathsf S^c_I$
onto a standard Borel parent space. It must satisfy five conditions.
\begin{enumerate}
\item[(L)] Laws. The parent belief and model laws and the parent evaluator
are functions of $R_I$ on $\Omega_I$.
\item[(D)] Dynamics. For a deterministic update, \cref{inh:fiber-test}
holds for all $x,x'\in\Omega_I$. For a flow, $R_I\circ\Phi_t=\Phi^c_t\circ R_I$
on $\Omega_I$ for every $t\geq0$ in the forward existence interval, with a
parent flow $\Phi^c$ that is well posed, with unique solutions, on
$R_I(\Omega_I)$ in a declared chart. For stochastic dynamics, a single
parent kernel reproduces the retained process for every initial law in
$\mathcal P_I$: strong lumpability when $\mathcal P_I$ contains all laws,
and, for example, the intertwining of \cref{inh:intertwining} when
$\mathcal P_I=\{\nu\Lambda\}$.
\item[(B)] Responses. For every $a\in A_I$ and every admitted receiver action
$u$ of an adaptive protocol, the joint kernel of the retained response and
the next retained state factors through $R_I$ on $\Omega_I$. When
$\mathcal P_I=\{\nu\Lambda\}$ is restricted, this is replaced by output
intertwining: with $T_{a,u,o}(x,dy)=\Pr(D_a(O)\in o,X'\in dy\mid X=x)$ for
measurable response sets $o$, require $\Lambda T_{a,u,o}=T^c_{a,u,o}\Lambda$
with $T^c_{a,u,o}=\Lambda T_{a,u,o}\Phi_{R_I}$.
\item[(C)] Channels. Every channel and evaluator read by a protocol in $A_I$
or by the parent update factors through $R_I$, including its dependence on
the configuration.
\item[(G)] Geometry, when declared. On an admitted regular family $R_I$
restricts to a smooth $G$-equivariant map independent of the context, and
the fine minimizing class is a coarse minimizing class, that is,
\cref{geo:connection-retention} holds; when equality of the full affine
classes is required, also $\ker J=\ker\bar J$
(\cref{geo:retained-connection}). The accompanying contraction of the
minimal tensor, with its exact defect, is \cref{geo:minimal-contraction}.
\end{enumerate}
This is the sense of inheritance in clause (F3) of
\cref{inh:recursive-formation}. An empty protocol class makes (B) vacuous,
and an unrestricted one forces the retention to be injective on exposed
child laws. Dynamical, response, and connection inheritance are separate
conditions. Output intertwining reduces to the factorization in (B) when
$\mathcal P_I$ contains all laws, and the induction in the proof of
\cref{inh:intertwining} goes through with $TP_z$ replaced by
$T_{a,u,o}P_z$, so it gives the predictive joint law of responses and
retained states and keeps the filter equal to $\Lambda(Z_n,\cdot)$ after
every observed response. The plain factorization in (B) would fail for
every response of a chain that is intertwined but not strongly lumpable,
such as the example after \cref{inh:intertwining}, because its second
marginal is $x\mapsto T(x,R_I^{-1}(\cdot))$.

This notion differs from a Markov-blanket hierarchy
\citep{KirchhoffEtAl2018,PalaciosEtAl2020,FristonEtAl2021}. A Markov blanket
is a conditional-independence property of a declared joint law over the
states of a system. A parent here is a retention that intertwines the given
dynamics and the declared responses; no joint law over configurations is
assumed or derived, and conditional independence enters only where a
separator is eliminated as in \cref{inh:separator-gluing}.

\definitionheading{Type closure}{inh:type-closed}
Fix a constitutive class $\mathcal C$ of agent dynamics, for example Fisher
gradient flows of reciprocal law potentials with normalized, possibly
law-dependent channels and evaluators. An inherited parent is type-closed
for $\mathcal C$ if its descended update, channels, and responses again
belong to $\mathcal C$, with the parent laws as agent laws and their Fisher
metric as the metric.

Type closure is what allows a result proved for primitive agents to be
applied to parents without change, and hence by induction at every level.
For the gradient class, \cref{inh:gradient-inheritance} gives a sufficient
test: the retention is horizontally conformal from the fine Fisher metric to
the Fisher metric of the parent laws, and the fine potential satisfies
\cref{inh:horizontal-criterion} with $a\equiv1$ and a reciprocal potential
of the parent laws. With a fiber-constant $a=\bar a\circ R\neq1$ the parent
is a gradient flow for $g_c/\bar a$, not for the parent Fisher metric.
Conformality is not necessary. The descended field $B^c$ is a
$g_c$-gradient exactly when the one-form $g_c(B^c,\cdot)$ is exact, which on
a simply connected chart is its closedness; type closure then requires in
addition that the potential belong to the declared class.

\openproblemheading{Type-closed parents}{inh:type-closure-open}
Determine which of the exact parents of the later chapters are type-closed
for the class of Fisher gradient flows of reciprocal law potentials. The
additive common-precision Wishart parent of \cref{ch:conditional-formation}
is answered by \cref{form:parent}, which identifies the descended field with
the unit-Fisher flow of the parent reciprocal potential by direct
computation. The conformality route does not reach it: on the
common-precision family the horizontal space of the sum retention forces the
conformal factor one in the precision directions, and the shape direction
would then need $\psi_2'(A_I)\sum_{i\in I}1/\psi_2'(\alpha_i)=1$, where
$\psi_2'(\nu)=\psi'(\nu)+\psi'(\nu-\tfrac12)$ is the shape Fisher
information of the bivariate Wishart law; for a block of at least two
agents the sum is not sufficient for the shapes, and the product is
strictly below one (it equals $0.480$ for two shapes equal to $1$). The
remaining parents of \cref{ch:smooth-closure} to
\cref{ch:collective-local-parent} are open.

\section{A finite two-level construction}

Let a complete configuration be $x=(a,b,c,r)\in\{0,1\}^4$. The coordinates
are labels of internal configurations, with $r$ an internal binary register
and $\rho(x)=r$ its readout. To attach actual agent laws, choose
\begin{equation}
 q_a=\operatorname{Ber}\left(\frac{1+2a}{4}\right),\qquad
 s_a=\operatorname{Ber}\left(\frac{2+a}{4}\right),\qquad
 \zeta_a=q_a\otimes s_a.
 \label{inh:finite-laws}
\end{equation}
Thus both sectors vary, with tied parameters in this supplied family;
$b,c$ are additional internal marks. A model presentation $m\in\{0,1\}$
has evaluator
\begin{equation}
 K(m,\{1\})=\frac14+\frac m2.
\end{equation}
The corresponding predictive law is
$\nu_a=s_aK=\operatorname{Ber}(1/2+a/8)$, which distinguishes the two
values of $a$.

At each step draw independent bits
$\xi\sim\operatorname{Ber}(1/4)$,
$\eta\sim\operatorname{Ber}(1/3)$, and
$\upsilon\sim\operatorname{Ber}(1/5)$, independently of preceding
draws, and set
\begin{equation}
 a'=a\mathbin\oplus\xi,\qquad
 b'=b\mathbin\oplus a\mathbin\oplus\eta,\qquad
 c'=c\mathbin\oplus b\mathbin\oplus\upsilon.
 \label{inh:finite-dynamics}
\end{equation}
Independently of these bit draws, conditional on the current configuration,
give $r'$ law $\nu_a$. This supplies a transition of an internal register;
the evaluator alone did not force this choice. The symbol $\oplus$ denotes
addition modulo two. The dynamics are coupled:
$a$ affects the next $b$, and $b$ affects the next $c$.

Retain $R_{10}(a,b,c,r)=(a,b,r)$ and then $R_{21}(a,b,r)=(a,r)$. For
$f_p(0)=1-p$ and $f_p(1)=p$, the first parent kernel is
\begin{equation}
 \begin{aligned}
 T^1((a,b,r),(a',b',r'))
 ={}&f_{1/4}(a'\mathbin\oplus a)\nu_a(\{r'\})\\
 &{}\cdot f_{1/3}(b'\mathbin\oplus b\mathbin\oplus a).
 \end{aligned}
 \label{inh:first-finite-parent}
\end{equation}
Summing the fine transition over $c'$ gives this expression, independently
of $c$. Summing it over $b'$ gives
$T^2((a,r),(a',r'))=f_{1/4}(a'\mathbin\oplus a)\nu_a(\{r'\})$,
independently of $b$.
Both reductions satisfy \cref{inh:finite-lumpability}. Their configuration
spaces have sizes sixteen, eight, and four, and both parents retain the laws
in \cref{inh:finite-laws}, their evaluator, and the register. Its readout
descends by the identity on record values. The first parent additionally
retains $b$.

Both parents reproduce the entire finite record history as well as the
retained configuration history. Indeed, write the fine path probability
as its initial mass times its transition factors. Summing over hidden
coordinates gives the parent path law by \cref{inh:finite-lumpability};
the record at each step is already a coordinate of that path. The next
record has probability $1/2+a/8$ of being one, so the inherited response
is nonconstant. No independence of the initial coordinates is needed.

This construction has nontrivial reductions and coupled dynamics, but its
retentions were specified in advance. It proves exact two-level
inheritance, without proving that the dynamics choose either reduction.
It is exact lumping of a drive--response chain: the fine configuration
carries one law pair and internal marks, so neither level groups two or
more agents, and each retained block is autonomous because the update of
$a$ reads only $a$ and the update of $(a,b)$ reads only $(a,b)$. For any
coordinate projection, \cref{inh:lumpability} is precisely such autonomy
of the retained block in law. In the reciprocal gradient dynamics of
\cref{ch:dynamics}, \cref{inh:coordinate-projection} shows that this
one-way mechanism is unavailable between energetically coupled agents.

\section{Averaged laws under smooth peer dynamics}

The arithmetic mixture of child laws is a normalized candidate parent law,
the linear opinion pool \citep{GenestZidek1986}, which minimizes the
average divergence $\sum_jw_j\KL(p_j\Vert q)$ \citep{Abbas2009}.
It need not determine either its own derivative or its boundary responses.
The following example establishes both failures within the smooth peer-KL
family of \cref{ch:dynamics}, with identical agents, a coupled evaluator,
and equal Fisher mobilities. It concerns this specified retention; it is
not an obstruction to every possible parent state.

\constructionheading{A coupled binary model with law-dependent channels}{inh:smooth-binary-model}
Use counting measure on $k,m\in\{-1,1\}$ and the strictly positive laws
\begin{equation}
 q_x(k)=\frac{1+kx}{2},\qquad s_y(m)=\frac{1+my}{2},
 \qquad x,y\in(-1,1).
 \label{inh:smooth-binary-laws}
\end{equation}
Choose $J=1/2$, reference $r(m)=1/2$, and normalized evaluator
$G_J(k\mid m)=e^{Jkm}/(2\cosh J)$. For the declared product recognition
family, the local variational term is exactly
\begin{equation}
 \begin{aligned}
 E_0(x,y)
 &=\sum_{k,m}q_x(k)s_y(m)
   \log\frac{q_x(k)s_y(m)}{r(m)G_J(k\mid m)}\\
 &=\psi(x)+\psi(y)-Jxy+\log\cosh J,\\
 \psi(x)&=\frac{1+x}{2}\log(1+x)
          +\frac{1-x}{2}\log(1-x).
 \end{aligned}
 \label{inh:smooth-local-energy}
\end{equation}
All expectations are finite. The mixed derivative is $-J$, so changing
either sector changes the force on the other.

In either sector, let a sender with parameter $x$ use the normalized kernel
\begin{equation}
 K_x(\ell\mid k)=\frac{1+\vartheta(x)k\ell}{2},\qquad
 \vartheta(x)=\frac{1+x^2}{4},\qquad
 q_xK_x=q_{f(x)},\quad f(x)=\frac{x+x^3}{4}.
 \label{inh:smooth-channel}
\end{equation}
Every entry exceeds $1/4$. The channel transmits the sampled sender bit,
and its transmission coefficient $\vartheta$ depends nontrivially on the
sender law. For passive frame covariance, present each alphabet as the atoms
$\{\pm u_i\}$ of a one-dimensional fiber, with $u_i\ne0$, and define
$\mathcal K_{ij}(u_jk,\{u_i\ell\})=K_{x_j}(\ell\mid k)$.
Under $g_i\in\mathrm{GL}(1,\mathbb R)$, push the laws forward and replace
$u_i$ by $g_i u_i$. The reference parameters remain fixed, and
\begin{equation}
 \mathcal K'_{ij}(g_jv,g_iA)=\mathcal K_{ij}(v,A).
 \label{inh:smooth-channel-covariance}
\end{equation}
The model sector, with its own representation of the same group, gives the same construction for $s$; pushing
$G_J$ through its model and belief frames preserves its normalization.
These frames are per-agent coordinates on each agent's fiber; the identity
is a passive re-expression and constrains nothing in the dynamics.
For two identical agents, set
\begin{equation}
 \begin{aligned}
 d(x,z)&=D_{\mathrm{KL}}(q_x\Vert q_z),\\
 P(x_1,x_2)&=d(x_1,f(x_2))+d(x_2,f(x_1)),\\
 E&=\sum_{i=1}^2 E_0(x_i,y_i)
       +P(x_1,x_2)+P(y_1,y_2).
 \end{aligned}
 \label{inh:smooth-peer-energy}
\end{equation}
Each source row contains only the other agent, so its normalized weight is
one and its row-relative-entropy term vanishes. This is a specialization
of the declared peer potential. The local term is the displayed KL
expectation; no further identification of the entire law-dependent
potential with a fixed-joint-law ELBO is assumed.
The product Fisher metric and its negative gradient flow are
\begin{equation}
 g=\sum_i\left(\frac{dx_i^2}{1-x_i^2}
                 +\frac{dy_i^2}{1-y_i^2}\right),\qquad
 \dot x_i=-(1-x_i^2)\partial_{x_i}E,\quad
 \dot y_i=-(1-y_i^2)\partial_{y_i}E.
 \label{inh:smooth-fisher-flow}
\end{equation}
Both mobilities are one. The vector field is smooth on $(-1,1)^4$.

\propositionheading{The mixture retention fails the dynamical inheritance test}{inh:smooth-mixture-failure}
For the flow in \cref{inh:smooth-fisher-flow}, the retention
\begin{equation}
 R(x_1,y_1,x_2,y_2)
 =\left(\frac{q_{x_1}+q_{x_2}}2,
         \frac{s_{y_1}+s_{y_2}}2\right)
 =(q_{\bar x},s_{\bar y})
 \label{inh:smooth-mixture-retention}
\end{equation}
does not support an autonomous parent vector field on the whole interior.
The same failure holds if the averaged joint law
$\frac12\sum_iq_{x_i}\otimes s_{y_i}$ is also retained.

\begin{proof}
For $0<|a|<1$ and $0<|c|<1$, take
$A=(0,c,0,c)$ and $B=(a,c,-a,c)$ in the coordinate order of
\cref{inh:smooth-mixture-retention}. Both retain $(q_0,s_c)$ and the
same averaged joint law $q_0\otimes s_c$. These are equal full binary
probability laws, not merely equal estimates of their means.
The identities $f(-x)=-f(x)$ and $d(-x,-z)=d(x,z)$ imply
\begin{equation}
 P(x_1,x_2)=P(x_2,x_1)=P(-x_1,-x_2),\qquad
 \partial_1P(a,-a)=-\partial_2P(a,-a),\quad
 \nabla P(0,0)=0.
\end{equation}
At $B$ the two Fisher prefactors are equal. Consequently the peer forces
and the two $\operatorname{atanh}$ terms cancel in $\dot{\bar x}$, giving
\begin{equation}
 \dot{\bar x}(A)=Jc,\qquad
 \dot{\bar x}(B)=Jc(1-a^2).
 \label{inh:smooth-mixture-defect}
\end{equation}
Both model-sector derivatives can also be written exactly. Define
\begin{equation}
 h(c)=2\operatorname{atanh}c-\operatorname{atanh}f(c)
       +f'(c)\frac{f(c)-c}{1-f(c)^2}.
\end{equation}
Indeed, direct differentiation gives
$\partial_xd(x,z)=\operatorname{atanh}x-\operatorname{atanh}z$ and
$\partial_zd(x,z)=(z-x)/(1-z^2)$. Thus
\begin{equation}
 \begin{aligned}
 (\dot y_1,\dot y_2)(A)&=-(1-c^2)(h(c),h(c)),\\
 (\dot y_1,\dot y_2)(B)&=-(1-c^2)(h(c)-Ja,h(c)+Ja),\\
 \dot{\bar y}(A)&=\dot{\bar y}(B)=-(1-c^2)h(c).
 \end{aligned}
\end{equation}
The retained belief derivatives nevertheless differ by $-Jca^2\ne0$.
For $a=1/2,c=1/4$, they are $1/8$ and $3/32$.
Hence $DR_x B(x)$ is not constant on a retention fiber, contradicting the
necessary condition \cref{inh:vector-square}. The identical retained
joint laws establish the additional assertion.
\end{proof}

The parameter $a$ may be arbitrarily small. A neighborhood of the
synchronized configurations therefore does not make this retention exact.
Unlike \cref{inh:role-loss}, the example uses identical dynamical roles.

\propositionheading{The same retention loses a boundary response}{inh:smooth-boundary-failure}
Choose a child uniformly, draw its belief bit, and transmit that bit
through \cref{inh:smooth-channel} to an external receiver. This declared
protocol has a normalized binary response law with mean
$b(x_1,x_2)=(f(x_1)+f(x_2))/2$. For any $m,a$ with
$|m|+|a|<1$, compare belief parameters $(m,m)$ and $(m+a,m-a)$,
keeping both model parameters equal to $c$. The retained marginal and
averaged joint laws agree, but
\begin{equation}
 b(m+a,m-a)-b(m,m)=\frac34ma^2.
 \label{inh:smooth-boundary-defect}
\end{equation}
For $ma\ne0$, their full binary response laws differ.

\begin{proof}
Expand $f(x)=(x+x^3)/4$ and use
$((m+a)^3+(m-a)^3)/2=m^3+3ma^2$.
At $m=a=1/4$, the two response means are $20/256$ and $17/256$.
Thus the response kernel is not constant on a retention fiber, violating
\cref{inh:boundary-square}.
\end{proof}

\propositionheading{Interior synchronization cannot first occur at finite time}{inh:smooth-no-finite-entry}
For \cref{inh:smooth-fisher-flow}, the diagonal
$\Delta=\{x_1=x_2,\ y_1=y_2\}$ is invariant. A solution starting outside
$\Delta$ cannot enter it at a finite time while remaining in the interior.

\begin{proof}
Agent exchange preserves the potential and metric, hence the vector field.
On its fixed set $\Delta$, the corresponding component derivatives agree,
which proves invariance. The smooth field has unique local solutions in
both time directions. If an interior solution meets $\Delta$, uniqueness
from that meeting point identifies it locally, backward as well as
forward, with the solution in $\Delta$. Along the original existence
interval, the set of times at which it lies in $\Delta$ is therefore open;
it is also closed by continuity. Connectedness of that interval makes it
the entire interval, contradicting the initial condition.
\end{proof}

The argument is the invariance of fixed-point subspaces of equivariant
vector fields \citep{GolubitskyStewartSchaeffer1988} combined with backward
uniqueness.

This result permits asymptotic synchronization, state-dependent threshold
selection, and exact quotients that do not require synchronization. It
does not prove convergence. An alternative retained state or a restricted
invariant family must be tested against both inheritance conditions;
the calculations above do not exclude such constructions.


\section{Memory and formation}

When a proposed retention fails, one exact option is to retain more of the
configuration. For a declared Markov process, with transition kernels that
may depend on time, another is to retain the conditional law of the current
configuration given the accessible retained history. That conditional law
determines future predictions. The Markov hypothesis is needed: for
independent fair bits $X_0,X_1$, $X_{n+1}=X_{n-1}$ for $n\geq1$, and
$R=\id$, the histories $(0,0)$ and $(1,0)$ have the same terminal
conditional law $\delta_0$ of $X_1$ but different next records. In a finite
model the conditional law is a probability vector with an exact recursive
update; on a countable retained alphabet the same formulas hold with
measures in place of vectors at every history of positive probability, and
on an uncountable one the update is a jointly measurable disintegration
defined almost surely (supplement). Its attainable values need not form a
finite set. Introducing such a filter presupposes a process and a
preparation law; it is not compulsory for deterministic coupled agents.

The conditional law is not a configuration retention in the sense of
\cref{inh:two-retentions}. It is a function of the retained history, and it
evolves as a Markov process on the simplex of laws on $\mathsf S$, an
information state \citep{Astrom1965}. In the finite case, with transition
matrix $T$ and the diagonal selector $P_z$ of $R^{-1}\{z\}$ as in the proof
of \cref{inh:intertwining}, it moves from $\widehat p$ to
$\widehat pTP_z/(\widehat pTP_z\mathbf 1)$ with probability
$\widehat pTP_z\mathbf 1$. It compresses only when noise or
nonobservability prevents the history from determining the fine state. For
deterministic dynamics observed through $R$, the filter is supported on the
image under $F^n$ of the initial configurations consistent with the history.
That consistent set decreases to the observability class of
\cref{inh:canonical-parent} with interface $R$, but its image need not
shrink. On a finite space, once the consistent set is a point the filter is
the point mass at the fine state; on infinite spaces the filter need not
concentrate even when the observability class is a point, as for the
doubling map $F(x)=2x\bmod1$ with $R(x)=1_{\{x\geq1/2\}}$ and a uniform
prior, where the filter on $X_n$ is uniform on an interval of length $1/2$
at every $n$. The intertwined handoffs of
\cref{inh:intertwining} are the case in which the filter is a function of the
current retained value.

\theoremheading{Information lost by a retained history}{inh:causal-memory}
Let $\mathbb P$ be a probability law on
$\mathsf Z_0\times\cdots\times\mathsf Z_{n-1}$, where each record space
is standard Borel and the horizon $n$ is finite. Set $H_t=(Z_0,\ldots,Z_{t-1})$,
with $H_0$ a point. For measurable maps into standard Borel statistic
spaces, let $M_t=f_t(H_t)$. Write
$P_t^M(\cdot\mid m)$ for a regular conditional law of $Z_t$ given
$M_t=m$. Let $\widetilde T_t$ be any normalized kernel from the
statistic space to $\mathsf Z_t$, and let $\widetilde{\mathbb P}$ be
the causal path law whose next-record kernel is
$\widetilde T_t(f_t(h),\cdot)$ at history $h$. Then, in $[0,\infty]$,
\begin{equation}
 \begin{aligned}
 \KL(\mathbb P\Vert\widetilde{\mathbb P})
 &=\sum_{t=0}^{n-1}I_{\mathbb P}(Z_t;H_t\mid M_t)\\
 &\quad+\sum_{t=0}^{n-1}\E_{\mathbb P}
 \KL\bigl(P_t^M(\cdot\mid M_t)
          \Vert\widetilde T_t(M_t,\cdot)\bigr).
 \end{aligned}
 \label{inh:causal-memory-identity}
\end{equation}
The first sum is the minimum path divergence among these predictors,
attained by choosing $\widetilde T_t=P_t^M$. Its $t$-th term vanishes
exactly when the true next-record law has a version depending on $H_t$
only through $M_t$.
\begin{proof}
Let $P_t(\cdot\mid h)$ be the true conditional law given $H_t=h$.
The relative-entropy chain rule expresses the path divergence as
\begin{equation}
 \sum_{t=0}^{n-1}\E_{\mathbb P}
 \KL\bigl(P_t(\cdot\mid H_t)
          \Vert\widetilde T_t(M_t,\cdot)\bigr).
 \label{inh:causal-row-chain}
\end{equation}
To decompose its $t$-th summand, compare the true joint law of
$(H_t,Z_t)$ with the law obtained by keeping its $H_t$ marginal and
drawing $Z_t$ from $\widetilde T_t(M_t,\cdot)$. Disintegrate these
two laws by $(M_t,Z_t)$. The divergence of their marginals is
$\E_{\mathbb P}\KL(P_t^M(\cdot\mid M_t)
\Vert\widetilde T_t(M_t,\cdot))$.
In the second joint law the conditional history given $(M_t,Z_t)$
is the true conditional history given $M_t$ alone. The conditional
divergence term is consequently
$I_{\mathbb P}(Z_t;H_t\mid M_t)$.

The chain rule on standard Borel spaces is an equality of nonnegative
extended quantities. A singular marginal or conditional contributes
infinite divergence; no subtraction of infinite terms is used. Applying
it to each summand of \cref{inh:causal-row-chain} gives
\cref{inh:causal-memory-identity}. Choosing the true statistic-conditional
rows makes every term in the second sum zero. Conditional mutual
information vanishes exactly when the corresponding conditional joint
law is the product of its conditional marginals. Since $M_t$ is a
function of $H_t$, this is equivalent to
$P_t(\cdot\mid H_t)=P_t^M(\cdot\mid M_t)$ almost surely.
\end{proof}

The finite-window case is the Markov approximation of a record process
\citep[Thm.~7.3.1]{Gray1990}; the theorem permits any measurable history
statistic. Its two terms separate forgetting information from predicting
incorrectly with the information kept. Vanishing of the first sum is
predictive sufficiency of the statistic \citep{ShaliziCrutchfield2001}. For
a statistic of a coarse process it implies, by data processing, the
dynamical-independence criterion of \citet{BarnettSeth2023}, whose target is
the next coarse value rather than the next record; the converse fails, as
for records $Z_t=(A_t,B)$ with independent fair bits $A_t$, one hidden fair
bit $B$, and coarse value $A_t$. A statistic alone need not have
a compact recursive realization. To update from the retained state and
new record only, one additionally needs
$M_{t+1}=U_t(M_t,Z_t)$, or a sufficient enlargement of $M_t$; this is the
determinism of causal states in \citet{ShaliziCrutchfield2001}.
The theorem concerns one preparation and protocol. Sufficiency for
every admitted intervention requires the corresponding conditional-law
identity for each such protocol.

A handoff initializes the parent from the actual fine state or its declared
conditional law, including the source marks and memory read by future
interactions. \Cref{inh:intertwining} gives the class of declared
conditional laws for which the parent is memoryless. When evidence is shared
across levels, its likelihood belongs to the joint model once. Reusing a
child's record through a complete same-episode parent posterior requires the
appropriate conditional or cavity accounting. A subsequent predictive prior
is a different operation.

No formed parent in \cref{ch:fusion} to \cref{ch:collective-local-parent}
acts as a source for its children; the supplied hierarchy of
\cref{ch:curvature-hierarchy} declares top-down factors instead. A
parent law can enter a child's cost as a law-dependent source, but the exact
reverse channel exerts no force: if $\zeta_I$ is the pushforward of
$\bigotimes_{i\in I}\zeta_i$ under a retained statistic and
$\mathcal K_{iI}$ is the regular conditional law of child $i$ given that
statistic, then $\zeta_I\mathcal K_{iI}=\zeta_i$, the marginal of a
disintegration. Downward action therefore needs a declared reverse channel
that differs from the exact conditional, and its inheritance must be tested
with the moving-map form of \cref{inh:vector-square}; no such construction
is given here.

\section{Statistical partitions and their limits}

An independence partition is a property of a declared joint law. A tuple
of individual agent laws does not itself specify a cross-agent joint
law; taking their product would impose independence in the input.
The following constructions apply when such a joint law has been supplied.

\propositionheading{The finest independent partition}{inh:finest-independent}
Let $V$ be a finite nonempty index set and let $\mathbb P$ be a
probability law on $\prod_{i\in V}\mathsf Y_i$, with standard Borel
coordinate spaces. Call a partition $\pi$ independent when
\begin{equation}
 \mathbb P=\bigotimes_{B\in\pi}\mathbb P_B,
 \label{inh:independent-partition}
\end{equation}
where $\mathbb P_B$ is its marginal on block $B$.
There is a unique finest independent partition $\pi_*(\mathbb P)$.
Relabeling the coordinates relabels this partition, and separate
bimeasurable changes of their presentations preserve it.
\begin{proof}
The one-block partition is independent. Suppose $\pi$ and $\sigma$
are independent. Marginalizing the $\sigma$ factorization to a block
$B\in\pi$ gives
\begin{equation}
 \mathbb P_B
 =\bigotimes_{\substack{D\in\sigma\\B\cap D\ne\varnothing}}
       \mathbb P_{B\cap D}.
\end{equation}
Substituting these marginals into the $\pi$ factorization proves
independence of the common refinement consisting of the nonempty
intersections $B\cap D$. There are finitely many partitions of $V$.
Refining all independent ones therefore gives an independent partition
finer than every one of them, which is necessarily unique.
Relabelings and coordinatewise bimeasurable maps commute with taking
marginals and product laws. Applying the inverse transformation gives
the reverse implication, so the set of independent partitions and its
unique finest member transform as stated.
\end{proof}

This statistical target is the finest-independent clustering problem
\citep{Ryabko2017}. Its existence does not require Gaussian laws or
densities, but does not identify statistical independence with agency.
In particular, it makes no statement about autonomous dynamics or
future boundary responses.

\corollaryheading{Conditional independence partitions}{inh:conditional-partition}
Let the same coordinate tuple have a joint law with an already retained
standard Borel variable $C$. Define
\begin{equation}
 \mathcal I_C(\pi)=
 \int\KL\left(
 \mathbb P(dy_V\mid c)
 \middle\Vert
 \bigotimes_{B\in\pi}\mathbb P_B(dy_B\mid c)
 \right)\mathbb P_C(dc).
 \label{inh:conditional-partition-information}
\end{equation}
There is a unique finest partition for which $\mathcal I_C(\pi)=0$.
\begin{proof}
Zero divergence means equality of the conditional joint and product
laws for $\mathbb P_C$-almost every $c$. There are finitely many
partitions. Outside the union of their relevant null sets, the
common-refinement argument in \cref{inh:finest-independent} applies
pointwise to every conditionally independent partition. Their common
refinement is therefore conditionally independent and uniquely finest.
\end{proof}

Choosing $C=Y_V$ makes every partition conditionally independent while
retaining the whole tuple. The corollary does not choose the context
that a useful parent should retain.

\propositionheading{Weak coupling can destroy every independent split}{inh:independence-discontinuity}
For any even $N\geq4$, let $y_i\in\{-1,1\}$ and define the strictly
positive finite law
\begin{equation}
 \mathbb P_\epsilon(y)
 =Z_\epsilon^{-1}\exp\left(
 J\sum_{b=1}^{N/2}y_{2b-1}y_{2b}
 +\epsilon\sum_{b=1}^{N/2-1}y_{2b}y_{2b+1}
 \right),
 \qquad J>0,\quad\epsilon\geq0,
 \label{inh:ising-partition}
\end{equation}
where $Z_\epsilon$ is the sum of the numerator over the finite alphabet.
Then $\pi_*(\mathbb P_0)$ consists of the pairs $\{2b-1,2b\}$, while
$\pi_*(\mathbb P_\epsilon)=\{V\}$ for every $\epsilon>0$, although
$\mathbb P_\epsilon\to\mathbb P_0$ in total variation.

\begin{proof}
At $\epsilon=0$, distinct pairs are independent and each pair is
dependent: its odds ratio is $e^{4J}\ne1$. Thus the finest partition
consists of those pairs.
For $\epsilon>0$ the interaction path is connected. Any nontrivial
partition cuts an edge. Fix all coordinates except the endpoints of
one cut edge, whose coupling is $J_e\in\{J,\epsilon\}$. Positivity
makes every such conditional law well defined. Its odds ratio is
$e^{4J_e}\ne1$, since all terms involving just one endpoint cancel
in that ratio. Factorization across the proposed partition would make
these endpoints conditionally independent after the other coordinates
were fixed, forcing ratio one. No nontrivial independent partition
exists. Continuity of each of the finitely many probabilities gives the
convergence in total variation.
\end{proof}

Exact independence is therefore not a continuous soft-membership criterion.

\propositionheading{Symmetry constrains a deterministic hard selector}{inh:selector-symmetry}
Let all coordinate spaces equal $\mathsf Y$, and suppose a deterministic
partition selector $\Pi$ reads only a law on $\mathsf Y^V$. Assume
relabeling naturality: $\Pi(\sigma_\#\mathbb P)=\sigma\Pi(\mathbb P)$ for
every permutation $\sigma$ of $V$. Let
$H_{\mathbb P}=\{\sigma:\sigma_\#\mathbb P=\mathbb P\}$ be the stabilizer of
the law. Then $\Pi(\mathbb P)$ is invariant under $H_{\mathbb P}$. If
$H_{\mathbb P}$ is transitive, the blocks of $\Pi(\mathbb P)$ form a block
system of $H_{\mathbb P}$; if it is primitive, $\Pi(\mathbb P)$ is either the
singleton or the one-block partition. In particular this holds when
$\mathbb P$ is invariant under all permutations.
\begin{proof}
For $\sigma\in H_{\mathbb P}$, naturality gives
$\Pi(\mathbb P)=\Pi(\sigma_\#\mathbb P)=\sigma\Pi(\mathbb P)$. An invariant
partition of a transitive group is by definition a block system, and a
transitive group is primitive when its only block systems are the two
trivial partitions \citep{DixonMortimer1996}. The full symmetric group is
primitive. For $|V|\leq2$ only the trivial partitions exist. For
$|V|\geq3$, if an invariant partition had a block $B$ with
$2\leq|B|<|V|$, choose distinct $i,j\in B$ and $k\notin B$. The
transposition exchanging $j$ and $k$ sends $B$ to a distinct block that
still contains $i$, which is impossible for a partition.
\end{proof}

For a law on a four-cycle invariant under its dihedral group, the only
nontrivial block system is the diagonal pairing $\{\{1,3\},\{2,4\}\}$, so a
relabeling-natural law-only selector cannot output adjacent pairs; adjacent
pairing needs asymmetric preparation data. A distribution over partitions,
soft memberships, asymmetric realized histories, or additional geometric
inputs are outside this proposition. An information-based hierarchy can also
be retained before choosing a cut \citep{ChanEtAl2016}. None of these
possibilities supplies the required dynamics without a further construction.
The symmetry result does not exclude symmetry breaking from asymmetric
initial conditions.

The four results of this section need a cross-agent joint law, and the
dynamics of \cref{ch:dynamics} evolves a tuple of agent laws with product
recognition. The product of the agent laws has the singleton partition as its
finest independent partition. The flow itself produces a joint law only as
the pushforward $(\Phi_t)_\#\mu_0$ of a preparation ensemble $\mu_0$, a law
over configurations that describes the ensemble rather than one population;
for a Dirac preparation it is a product of point masses whatever the
coupling. For a nondegenerate ensemble coupling makes it dependent across
interacting agents: for
$\dot x=Ax$ with $A=\bigl(\begin{smallmatrix}-1-\lambda&\lambda\\
\lambda&-1-\lambda\end{smallmatrix}\bigr)$ and $x(0)\sim\mathcal N(0,I)$,
$\operatorname{Cov}(x_1(t),x_2(t))=\tfrac12(e^{-2t}-e^{-2(1+2\lambda)t})>0$
for $\lambda,t>0$. The theory also supplies content-level cross-agent joint
laws: the posterior $\Pi_o$ of the global model of \cref{sec:dyn-likelihood}
and the joint private-record posterior of \cref{app:interacting-memory}.
Both are admissible inputs for \cref{inh:finest-independent} and
\cref{inh:selector-symmetry}, and their finest partitions are determined by
the declared model or protocol rather than by the flow of agent laws. These results are therefore constraints
on admissible selectors rather than formation mechanisms; the selectors of
the later chapters are declared thresholds for this reason.

\section{Architectures whose selection cannot see the interaction}

Two exact statements constrain the membership order that some architectures
can produce. They concern the full flow; whether a construction then fails
the interaction-necessity clause (F5) of the formation standard depends in
addition on its null flows and on the class of selectors it declares, and
the remarks after each proposition say when it does.

\propositionheading{Scalar additive closure fixes the membership order}{inh:scalar-additive}
Let scalar parameters $\theta_i$ in an open interval $\mathsf J$ evolve by
$\dot\theta_i=G(\theta_i,g(t))$, where $g$ is a state common to all agents and
$G$ is continuous in its first argument. Suppose that for every pair of
agents the sum retention $\theta_i+\theta_j$ intertwines with the flow, that
is, $G(x,g)+G(y,g)$ depends on $x,y\in\mathsf J$ only through $x+y$. Then
$G(\theta,g)=a(g)\theta+c(g)$, and every difference $\theta_i-\theta_j$ and
every difference of count-weighted block centroids equals its initial value
times the common factor $m(t)=\exp\int_0^ta(g(\tau))d\tau$. Consider a
selector that joins two current blocks when their centroid distance first
reaches a fixed threshold from above, and suppose $m$ is nonincreasing and
the initial centroid distances have no ties. Then the sequence of selected
partitions is an initial segment of a sequence determined by the initial
data alone; the common state $g$ determines only the event times and the
length of the segment.
\begin{proof}
For fixed $g$ put $\phi=G(\cdot,g)$. If $x+y=x'+y'$ then
$\phi(x)+\phi(y)=\phi(x')+\phi(y')$; taking $x'=y'=(x+y)/2\in\mathsf J$ gives
Jensen's equation $\phi((x+y)/2)=(\phi(x)+\phi(y))/2$, whose continuous
solutions are affine. Hence $\frac{d}{dt}(\theta_i-\theta_j)=a(g)(\theta_i-\theta_j)$,
and a count-weighted centroid obeys the same affine law, so centroid
differences are also multiplied by $m(t)$. A pair of current blocks with
initial centroid distance $D$ is at distance $Dm(t)$ and first reaches the
threshold $\delta$ when $m(t)=\delta/D$. Since $m$ is common and
nonincreasing, events occur in increasing order of $D$. After a join at
level $D$, a pair involving the new block with initial centroid distance
$D'\leq D$ is already at or below the threshold and, $m$ being
nonincreasing, never reaches it from above; a pair with $D'>D$ joins later in
order of $D'$. All of these comparisons involve only initial distances.
Whether the event at level $D$ occurs depends on whether $m$ reaches
$\delta/D$, which is where $g$ enters.
\end{proof}

The same Jensen argument appears as \cref{form:scalar-closure} for the
Wishart family. The proposition says nothing about the null flows, whose
force need not be affine or share the factor $m$. Under scalar additive
closure the full-flow membership order is fixed by the initial data;
interaction necessity fails in addition when the null flow is scalar
additive with a common factor $m^0$ and the infima of $m$ and $m^0$ select
the same segment, for example when both tend to zero and histories are
compared on $[0,\infty)$. When the null factor is $m^0\equiv1$, or the two
infima lie between different event levels $\delta/D_k$, the partition
sequences differ on open sets although the full-flow order is a dendrogram
of the initial data; the calibration example after the formation standard is
of this kind.

Anisotropic linear laws escape this proposition. For vector parameters with
$\dot\theta_i=A(g)\theta_i+c(g)$, differences evolve by a common propagator
that need not preserve the ordering of distances: with
$A=\operatorname{diag}(-1,-3)$, initial differences $(1,0)$ and $(0,1.5)$
become approximately $(0.37,0)$ and $(0,0.075)$ at $t=1$, reversing which
pair is closer. Within the hypotheses of \cref{inh:common-clock} below this
cannot happen: on finite tables the centered log-ratios are the
natural-parameter differences, and they obey a common propagator equal to
$-2$ times the identity whatever the dimension and the coupling. Whether a
smooth reciprocal potential outside that class realizes an anisotropic,
interaction-dependent propagator with exact additive closure is open.

\propositionheading{Shared targets give a common clock}{inh:common-clock}
Work on finite positive tables as in \cref{thm:dyn-finite-flow}, with a
common belief alphabet $\mathsf K$. Suppose every evaluator is the same
model-independent law, $G_i(m,\cdot)=G_0$, and every belief channel from a
source $j\neq i$ is $\mathcal K^q_{ij}[\mathbf q](v,\cdot)=C[M]$ for all
inputs $v$, where $C$ is a twice continuously differentiable normalized law
of the additive moment $M=\sum_i\E_{q_i}\mathsf T$ of a fixed statistic
$\mathsf T:\mathsf K\to\R^p$. Suppose either (a) the belief source sets
exclude self, or (b) they include self with the self channel also equal to
$C[M]$. Then along the Fisher flow of \cref{eq:dyn-potential} (with dynamic
rows) or of its reduced potential \cref{eq:dyn-reduced},
\begin{equation}
 \frac{d}{dt}\log\frac{q_i(k)}{q_j(k)}
 =-2\log\frac{q_i(k)}{q_j(k)}+c_{ij}(t),
 \label{inh:common-clock-identity}
\end{equation}
with $c_{ij}$ independent of $k$. In case (a) the two nulls of
\cref{inh:null-flow} coincide, and along them the coefficient $-2$ is
replaced by $-1$ with the common target $G_0$. In case (b) the self costs
are nonzero and the nulls differ: along the label-preserving null the
coefficient of $\log q_i$ in the force is $1+\beta^0_{ii}$, with the agent-
and time-dependent self weight of \cref{inh:null-self-row}, and along the
renormalized null it is $2$ with a forcing that depends on the agent through
the frozen peer moment; neither has a common clock in general.
\begin{proof}
With a model-independent evaluator the local term splits as
$\KL(s_i\Vert r_i)+\KL(q_i\Vert G_0)$, whose belief derivative is
$\log q_i-\log G_0$ plus a constant. In both cases every cost in the row of
$i$ is $\KL(q_i\Vert C[M])$ and contributes $\beta_{ij}(\log q_i-\log C[M])$
as receiver, and the row sums to one. As a source, $q_i$ enters every cost
only through $M$, so it contributes
$-\sum_{l}\sum_{j}\beta_{lj}\sum_u\frac{q_l(u)}{C[M](u)}
\frac{\partial C[M](u)}{\partial M}\cdot\mathsf T(v)$, the sum running over
the source sets, which is the same function of $v$ for every $i$; with Gibbs
rows the envelope identity \cref{eq:dyn-envelope} gives the same expression
with $\beta^*$. Hence $h_i=2\log q_i+V_t$ with $V_t$ common to all agents, up
to constants in $v$. The flow \cref{eq:dyn-fisher-flow} gives
$\frac{d}{dt}\log q_i=-(h_i-\E_{q_i}h_i)$, and subtracting the equations for
$i$ and $j$ proves \cref{inh:common-clock-identity}. In case (a) deleting
the peer costs leaves no row term under either null, so
$h_i=\log q_i-\log G_0$. In case (b) the surviving self cost is
$\KL(q_i\Vert C[M_i^0])$ with $M_i^0=\E_{q_i}\mathsf T+M^0_{-i}$ and the
initial peer moment $M^0_{-i}$; it enters with weight $\beta^0_{ii}$ or one,
and $M^0_{-i}$ differs between agents.
\end{proof}

\Cref{form:mean-field} is the instance of case (b) on the Wishart family,
where only the shape coordinates keep a common clock. In case (a) the proof
gives more than the identity.

\corollaryheading{The full flow is a common tilt of the null flow at doubled time}{inh:common-tilt}
In case (a) of \cref{inh:common-clock}, let $q^0_i$ be the null flow from the
same initial laws. There are positive functions $\omega_t$ on $\mathsf K$,
the same for every agent, such that
\begin{equation}
 q_i(t)=\frac{\omega_tq_i^0(2t)}{\E_{q_i^0(2t)}\omega_t}
 \qquad\text{for every }i\text{ and }t\geq0.
 \label{inh:common-tilt-identity}
\end{equation}
\begin{proof}
From $h_i=2\log q_i+V_t$, $\log q_i(t)=e^{-2t}\log q_i(0)+W_t+c_i(t)$ with
$W_t=-\int_0^te^{-2(t-s)}V_sds$ common and $c_i$ constant on $\mathsf K$.
Along the null, $\log q_i^0(\tau)=e^{-\tau}\log q_i(0)+(1-e^{-\tau})\log G_0
+c_i'(\tau)$. At $\tau=2t$ the difference of the two logarithms is
$W_t-(1-e^{-2t})\log G_0$ up to constants on $\mathsf K$, which does not
depend on $i$; take $\omega_t$ its exponential.
\end{proof}

A selector that is invariant under common tilts and under monotone changes
of time, such as one that reads only pairwise log-ratios modulo constants,
therefore produces along the full flow the partition sequence it produces
along the null flow, on every horizon, and fails (F5) in this class. A
selector that is not tilt-invariant, such as a threshold on the divergence
between two agent laws, can give different partition sequences although
every relative coordinate evolves as along the null flow. This is one reason
why the formation standard fixes a selector family before the flows are
analyzed and asks for a selector-free observable. Interaction-sensitive
selection under the reciprocal potential, read by tilt-invariant selectors,
therefore requires that at least one hypothesis of \cref{inh:common-clock}
fail: heterogeneous or model-dependent evaluators, self sources with nonzero
self costs as in case (b), input- or source-dependent channels, law
dependence beyond one additive moment (a shared target that is the
normalized geometric mean of all agent laws is the same law for every
receiver and is not of this form), or receiver-dependent targets such as
selective channels with nonconstant Gibbs rows. Heterogeneous evaluators
alone can give (F5) with a membership order supplied by the references, as
in the calibration example below.

\section{The formation standard}

The thesis of the theory is that interactions between agents can produce
further agents. A definition of formation must therefore exclude three ways
in which the selection and inheritance clauses can be met without any such
production. A parent that retains every child state compresses nothing:
four agents with laws $\mathcal N(x_i,1)$ under symmetric peer-KL consensus,
a threshold selector, and the identity retention select two pairs and then
their union with exact inheritance at every stage. Agents that do not
interact can be grouped by any rule that reads their initial data and a
common relaxation clock. A second stage that forms from every handoff
whatever the first stage produced is a sequence of supplied scales rather
than recursion. The standard below makes each exclusion a clause, and it
fixes before the dynamics every object that a construction could otherwise
choose after seeing its own flows: the controls, the topology of the
admitted class, the protocol and probe classes, the selector and its
constants, the horizon, and the counterfactual class.

\definitionheading{Null flows and the frozen-environment control}{inh:null-flow}
Let a construction declare its objective and its flow $\Phi$, with
references, evaluators, source sets and priors, temperatures, records, a
selector family, and an admitted class of initial configurations. A cost of
receiver $i$ is a peer cost when it belongs to an edge $j\to i$ with
$j\neq i$ whose compared law $\nu_{ij}$ depends, on the admitted class, on
the law of some agent other than $i$. A cost carrying a peer label whose
compared law depends on no other agent's law, such as one with a
constant-output channel, is a self term and is kept by every control.
\begin{enumerate}
\item[(N1)] The label-preserving null $\Phi^0_{\mathrm{lp}}$ deletes every
peer cost, that is, sets it to zero, while the peer labels, the source
priors, the temperatures, and all other data are unchanged. Every
law-dependent channel and evaluator that remains is evaluated with the laws
of the other agents held at their initial values. With Gibbs rows, a
receiver with self prior $a$ and self cost $D_{ii}$ then has self weight
\begin{equation}
 \beta^0_{ii}=\frac{ae^{-D_{ii}/T}}{ae^{-D_{ii}/T}+1-a},
 \label{inh:null-self-row}
\end{equation}
which is below one whenever $a<1$, and the same form holds in the model
sector.
\item[(N2)] The renormalized null $\Phi^0_{\mathrm{rn}}$ is the same, except
that the labels of the deleted edges are removed from the source sets. A
receiver whose source set contains itself then has self weight one, and a
receiver with no remaining source has no row term.
\item[(N3)] The frozen-environment control $\Phi^{\mathrm{fr}}$ deletes
nothing. In every cost, row, channel, and evaluator that enters the force on
agent $i$, each law of an agent $j\neq i$ is held at its initial value.
\end{enumerate}
The two nulls coincide when no receiver has a self source whose cost is not
identically zero. In the per-agent picture the same rules apply to each
$\mathcal F_i$ of \cref{prop:dyn-per-agent}, for the dynamics of
\cref{def:dyn-per-agent-flow} and for its Gibbs-reduced version in which the
rows are minimized exactly at unit temperature at every state; the self
costs are $\KL(s_i\Vert r_i)$ and $\int\KL(q_i\Vert G_i(m,\cdot))s_i(dm)$,
and (N1) keeps the echo labels. For the global free energy of
\cref{sec:dyn-likelihood}, write the coordinates of the model as the agent
blocks $y_i$ and the non-agent coordinates $\nu$, such as structural latents,
block latents, and source labels not attached to a single agent. Both nulls
marginalize every record whose likelihood factor depends on the contents of
two or more agents and keep single-agent factors and records. The
label-preserving null replaces $P_0$ by the product over $i$ of the marginal
of $P_0$ on $(y_i,\nu)$: each agent receives a private copy $\nu^{(i)}$ of the
non-agent coordinates with their prior, so latent source labels and
structural latents are kept at their priors and are read only by their own
agent, and a single-agent factor that also reads $\nu$ reads that agent's
copy. The renormalized null integrates the private copies out and replaces
$P_0$ by the product of its marginals on the agent blocks alone. The two
nulls coincide when the model has no non-agent coordinate and differ in
general in their transients. The frozen-environment control gives every
agent $i$ a private copy of the non-agent blocks, updated by the rule that
updates them along the full flow, and holds the law of every agent $j\neq i$
at its initial value in each factor that enters the force on $Q_i$ or on its
copy; without non-agent coordinates this is $Q_{-i}$ held at its initial
value in \cref{eq:dyn-mean-field-field}. A selector that reads a non-agent
block, such as a structural recognition factor, reads along every control the
block minimizer that the original objective gives at the control's agent
laws, as a row-reading selector does. Every control is run in the
recognition family and with the update schedule of the full flow.
\Cref{ch:global-formation} states a selection result for the global free
energy under these conventions.

Each control is the flow of the modified objective in the chart or invariant
family in which the construction declares $\Phi$, and the construction
verifies for it the Lipschitz and well-posedness hypotheses it verifies for
$\Phi$ (for example those of \cref{prop:dyn-chart-flow}). The control must
exist, with unique solutions, on the admitted class up to the declared
horizon; where it does not, (F5) is undefined there and the admitted class is
restricted. A construction whose objective lies outside these classes, such
as the nonsmooth interaction of \cref{ch:fusion}, declares its controls by
the same rules and states their well-posedness.

The selector is a declared measurable function of the configuration of agent
laws, the accessible records, and the declared interfaces, built from
constitutive data and the same function at every stage. Rows and costs enter
it only through the original constitutive formulas as functions of the laws,
and it is evaluated in the same way along $\Phi$ and along every control: a
selector that reads rows reads, along a control, the rows that the original
costs give at the control state, not the rows of the control's own
objective. After a handoff along a control, the control remains the fine
flow, and the selector is evaluated on the block statistics it reads,
computed from the fine states; inheritance is not required along a control,
and where retained parent laws are not defined along it the construction
extends the selector to fine states by a declared formula.

The membership history $H(x)$ of an initial configuration $x$ is the
sequence of selected partitions with their event times on the declared
horizon $[0,\tau]$, or on the maximal forward existence interval when no
finite horizon is declared. The partition sequence $P(x)$ is $H(x)$ with its
event times removed, two sequences being identified when one is obtained
from the other by exchanging consecutive events whose changes involve
disjoint sets of agents. The histories $H^0_{\mathrm{lp}}$,
$H^0_{\mathrm{rn}}$, $H^{\mathrm{fr}}$ and the sequences $P^0_{\mathrm{lp}}$,
$P^0_{\mathrm{rn}}$, $P^{\mathrm{fr}}$ are produced by the same selector
along the controls, each on its own maximal forward existence interval.

\lemmaheading{Agents that do not interact fail interaction necessity}{inh:no-interaction-null}
If no cost, row, channel, or evaluator that enters the force on any agent
depends on the law of another agent on the admitted class, then
$\Phi^0_{\mathrm{lp}}=\Phi^0_{\mathrm{rn}}=\Phi^{\mathrm{fr}}=\Phi$ there, and
$P=P^0_{\mathrm{lp}}=P^0_{\mathrm{rn}}$.
\begin{proof}
By the definition there is no peer cost, so the nulls delete nothing, and
holding the laws of other agents fixed changes no term of any force. The
selector is the same function along all four flows.
\end{proof}

The rule that peer costs are those with a compared law depending on another
agent is needed for this lemma. Deleting every cost with a peer label would
remove, for three Bernoulli agents with a common evaluator and a
constant-output channel from each peer, a term $\KL(q_i\Vert C)$ that couples
no agents but moves every agent's own target, and a relabeling-equivariant
selector reading the distance to the evaluator then gives different
partition sequences along the full flow and that deletion on an open set.
The frozen-environment control answers a different question from the nulls:
whether the selection already follows from each agent's response to the
initial laws of the others, without feedback.

\definitionheading{Counterfactual handoff class}{inh:counterfactual-class}
For a construction with two selection stages, let the first stage hand off
at a time that is deterministic or a stopping time, with its selected blocks
and retentions, and let the second-stage event and its block be measured from
the handoff. A counterfactual class $\mathcal C$ is declared with the
constitutive data, before the second stage is analyzed. It contains every
actual first-stage handoff state reachable from the admitted class and the
counterfactual handoff states obtained by applying the same retentions, at
the same handoff times, to the states of trajectories from a declared class
$\mathcal U$ of initial configurations containing the admitted class, along
which the first-stage event has not occurred. It excludes states outside the
domain of the selector, such as parents already within a threshold that the
selector requires to be reached from above. It carries the relative
topology of the parent configuration space and, for stochastic dynamics, the
declared preparation law of the counterfactual handoffs, under which they
have probability bounded below.

\definitionheading{Inherited recursive formation}{inh:recursive-formation}
An inherited recursive formation is a construction satisfying all of the
following clauses.
\begin{enumerate}
\item[(F1)] Admitted class. The initial configurations form a nonempty set of
initially unpartitioned configurations that is open in the chart topology of
the declared dynamics: the product of open simplices for finite tables, the
bounded log-density chart of \cref{def:dyn-chart-flow}, a declared Banach
chart, or the parameter manifold of a declared restricted-family flow, which
is then stated. This is the full grade. In the family grade the set is
instead relatively open in a declared invariant family defined by stated
constitutive constraints and not contained in finitely many orbits, and its
codimension is reported. The grade is always reported.
\item[(F2)] Selection. The dynamics selects nontrivial blocks with a selector
from a declared family $S_c$, $c\in\mathsf C$, that satisfies the readout
convention of \cref{inh:null-flow}, with a stated persistence criterion: an
invariant block relation, persistence on a finite window, or a declared path
event for stochastic dynamics. The family and the set $\mathsf C$ of its
constants are fixed before the flows are analyzed, and every real constant
ranges over a set with nonempty interior. Each handoff is a measurable
function of the state or of accessible history.
\item[(F3)] Inheritance. Each selected block carries an inherited parent
(\cref{inh:inherited-parent}) for the classes of \cref{inh:protocol-classes}.
\item[(F4)] Compression. For each handoff there are $b\in B_I$ and
$x\neq x'$ in the reachable post-handoff set $\Omega_I^{\mathrm{reach}}$ with
$R_I(x)=R_I(x')$ and $\operatorname{law}_x(b)\neq\operatorname{law}_{x'}(b)$.
Equivalently, $R_I$ is noninjective on $\Omega_I^{\mathrm{reach}}$ and erases
a distinction that is not gauge relative to $B_I$.
\item[(F5)] Interaction necessity. Against a null $\Phi^0$ of
\cref{inh:null-flow} and for a constant $c\in\mathsf C$, there is a nonempty
open subset $U$ of the admitted class such that, for every $x\in U$,
$P(x)\neq P^0(x)$ when both histories are taken on $[0,\infty)$, and $P(x)$
on the declared finite horizon $[0,\tau]$ is not an initial segment of
$P^0(x)$ when a finite horizon is declared. The clause holds when this is
true against both nulls and for every $c\in\mathsf C$.
\item[(F6)] Recursion. (F6a) The parents select a block containing at least
two first-stage parents under the inherited interactions, by the
first-stage selector family with the same constants and the same
constitutive data, without resetting their source data. (F6b) For the
declared counterfactual class $\mathcal C$, the second-stage event occurs
after the actual first-stage handoffs and does not occur from a set of
counterfactual handoffs with nonempty interior in $\mathcal C$; for
stochastic dynamics, the probability of the second-stage event by the
declared windows after the actual handoffs exceeds its value after the
counterfactual handoffs by a margin bounded below.
\end{enumerate}
A construction satisfying (F1), (F2), (F3), and (F6a) and shown to fail at
least one of (F4), (F5), and (F6b) is a two-stage selection with exact
inheritance; a construction with only one selection stage that satisfies
(F1) to (F3) is a single-stage selection with exact inheritance. A clause
that is neither established nor shown to fail leaves the classification
undetermined, and the construction is then not shown to be a recursive
formation. The clauses that fail or remain undetermined are always stated.

Clause (F5) is a control experiment: it asks whether the grouping would
occur with the interaction switched off. Its verdict names the null. When
the two nulls give different verdicts it is reported as split, and the
clause fails; a pass against the label-preserving null alone can come from
the reweighting \cref{inh:null-self-row} of a nonzero self cost, which
carries no information about the interaction. The verdict reports the subset
of $\mathsf C$ on which it holds, a selector-free discrepancy between $\Phi$
and the nulls where one exists (a limit point, a sign of a public coordinate,
a contraction factor), and the comparison of $P$ with $P^{\mathrm{fr}}$.
A difference only in event times, or only in the order of merges of
disjoint blocks, does not count, and the initial-segment comparison on a
finite horizon excludes a pure change of clock: under the common clock of
\cref{inh:common-clock} the full flow realizes on $[0,\tau]$ the null
sequence up to the null time $2\tau$, which a comparison of both histories on
$[0,\tau]$ would count as a difference. Persistence can mean an invariant
block relation in deterministic dynamics or a specified path event in
stochastic dynamics, and a finite-window criterion is stated as such.

Clause (F4) excludes the identity retention described above, and it is
always reported with two attributes: whether the separation by $b$ is
persistent or tends to zero along the post-handoff flow, and whether the
fibers of $R_I$ are orbits of a symmetry of the fine field and of the
protocols of $A_I$. Under (F3) the retained dynamics and every admitted
parent response are blind to the erased distinction, while (F4) asks that a
child-level probe declared before the handoff see it. The exact routes of
this chapter, symmetry quotients and horizontal energy structure, produce
such compressions; no mechanism is given here by which an interaction
creates an exact symmetry or decoupling on an open admitted class, and a
graded notion of inheritance is not developed.

In (F6), a counterfactual class that contained parents already within the
threshold of a from-above selector would make (F6b) hold for every such
construction, since those parents never unite; the domain exclusion in
\cref{inh:counterfactual-class} removes this case. If the second-stage event
occurs from every state of $\mathcal C$ outside a set with empty interior, or
for stochastic dynamics with probability tending to one uniformly over
$\mathcal C$, the stages form a supplied-scale sequence; a proof of the
second stage that uses only hypotheses on the handoff state that also hold
for counterfactual handoffs in $\mathcal C$ establishes such a sequence. A
margin in the other direction, where the first stage delays or suppresses
the second, is first-stage dependence and not recursion.

Two grades refine (F5). The strong grade holds when (F5) holds at every
point of the admitted class, for every $c\in\mathsf C$, against both nulls.
For the amplifying grade a construction declares a metric $d$ on agent laws,
for example the Hellinger metric, and for a partition $\pi$ of a
configuration $y$ puts $\varrho(y,\pi)$ equal to the largest within-block
$d$-diameter divided by the smallest $d$-distance between blocks. It is
amplifying with factor $\lambda<1$ if, at every event time $t_k$ of $\Phi$
with partition $\pi_k$ and for both nulls,
$\varrho(\Phi_{t_k}x,\pi_k)\leq\lambda\varrho(\Phi^0_{t_k}x,\pi_k)$ on an
open set of the admitted class. Both ratios are evaluated on the partition of
$\Phi$, so the grade is defined when $\Phi^0$ has no event; when within-block
diameters vanish at the events of $\Phi$ the grade holds trivially and is
reported as degenerate.

The standard does not require that a parent be type-closed in the sense of
\cref{inh:type-closed}, that it act on its children through a term of the
fine field absent before the handoff, or that it hold a generative model or
lower a free energy. Two clauses of (F6) therefore establish two stages with
inherited dynamics, not iteration to every depth or top-down efficacy;
whether these become clauses is part of \cref{inh:formation-open}.

Each construction of the later chapters is measured against this standard
with the following report: its objective, dynamics, and chart; the grade of
(F1); the selector family, its constant set, and the persistence criterion;
the rules of \cref{inh:protocol-classes}; the witness and the attributes of
(F4); the verdict of (F5) against each null over $\mathsf C$, with the
selector-free observable and the frozen-environment comparison; (F6a); the
counterfactual class and the verdict of (F6b); the grades; and the resulting
classification.

\propositionheading{An anchored linear consensus is a two-stage selection that fails (F6b)}{inh:anchored-consensus}
Take four agents with belief laws $\mathcal N(x_i,1)$, $x\in\R^4$, a trivial
model sector, evaluators $\mathcal N(c_i,1)$ with
$c=(0,0.6,1.3,1.95)$, identity channels including the self channel, and
fixed rows of weight $1/4$ on each of the four sources. The reciprocal
potential is
\begin{equation}
 \mathcal E=\sum_i\tfrac12(x_i-c_i)^2
 +\tfrac14\sum_i\sum_{j\neq i}\tfrac12(x_i-x_j)^2,
 \label{inh:anchored-potential}
\end{equation}
and the declared dynamics is its restricted-family flow
(\cref{def:dyn-restricted-flow}) $\dot x=-\nabla\mathcal E$, with unit Fisher
metric. The selector family joins two current blocks when the distance of
their count-weighted centroids reaches $\delta$ from above,
$\delta\in\mathsf C=[0.45,0.55]$; the retention of a block is its centroid
and count, all other coordinates kept; $A_I$ consists of the average of one
fresh sample from each member, and $B_I$ adds one fresh sample from a named
member through the identity channel the population uses. Then, on a
neighborhood $U$ of $c$ and for every $\delta\in\mathsf C$:
\begin{enumerate}
\item[(i)] From $x(0)=c$ every difference of coordinates or centroids equals
its initial value times $m(t)=(1+2e^{-3t})/3$. The full flow selects
$\{1,2\}$, then $\{3,4\}$, then $\{1,2,3,4\}$, at the times with
$m(t_k)=\delta/D_k$, $D_k=0.6,0.65,1.325$; every relation persists. The two
nulls coincide and give $x(t)=c+e^{-t}(x(0)-c)$, which selects nothing. The
frozen-environment control gives the factor $(1+4e^{-5t/2})/5$ from $x(0)=c$
and the same partition sequence as the full flow.
\item[(ii)] (F1) holds in the full grade for the declared restricted flow,
(F2), (F3), and (F6a) hold, (F4) holds with a transient separation that is
not of symmetry type, (F5) holds in the strong grade against both nulls with
selector-free observable the contraction of every difference by the factor
$1/3$, and $P^{\mathrm{fr}}=P$.
\item[(iii)] For every counterfactual class of \cref{inh:counterfactual-class}
the union of the two parents occurs from every counterfactual handoff, so
(F6b) fails and the construction is a two-stage selection with exact
inheritance.
\end{enumerate}
\begin{proof}
The flow is $\dot x_i=-(x_i-c_i)-\tfrac12\sum_{j\neq i}(x_i-x_j)$, the source
role contributing the same term as the receiver role. Since
$\sum_{j\neq i}(x_i-x_j)=4x_i-\sum_jx_j$, every difference obeys
$\dot\Delta=-3\Delta+(c_i-c_j)$, so it relaxes at rate three to one third of
the corresponding anchor difference, and from $x(0)=c$ it equals
$(c_i-c_j)m(t)$; count-weighted centroids are averages and obey the same
law. Since $m$ decreases from $1$ to $1/3$, a pair of current blocks with
initial centroid distance $D$ reaches $\delta$ exactly when $m=\delta/D$,
which happens for $D/3<\delta<D$. The initial distances are $0.6$ for
$\{1,2\}$, $0.65$ for $\{3,4\}$, at least $0.7$ for the other pairs, and after
the joins $1.0$ and $1.65$ from $\{1,2\}$ to the singletons and $1.325$
between the two parents. For $\delta\in\mathsf C$ the order is as stated, the
last event requires $\delta>1.325/3=53/120$, and $m$ decreasing makes every
joined relation invariant. Along both nulls the self cost vanishes and the
rows are fixed, so deleting the peer costs, with or without their labels,
leaves $\dot x=-(x-c)$; its differences move monotonically from their initial
values to the anchor differences, which are at least $0.6>\delta$. Along the
frozen-environment control,
$\dot x_i=-\tfrac52x_i+c_i+\tfrac12\sum_{j\neq i}x_j(0)$, and from $x(0)=c$
the differences equal $(c_i-c_j)(1+4e^{-5t/2})/5$, a decreasing factor with
limit $1/5<0.45/1.325$, which gives the same sequence. Every crossing is
transverse, since the factors have negative derivative, and at every event
the other current distances exceed $\delta$ strictly, uniformly for $\delta$
in the compact set $\mathsf C$; continuous dependence of the linear flows on
initial data therefore gives the same three sequences on a neighborhood $U$.

For (ii), $U$ is open in $\R^4$, the chart of the declared flow, and is
initially unpartitioned. Each centroid evolves by the averaged law, which
reads only centroids and counts, and the force of a block $I$ on an outsider
$k$ is $\tfrac12|I|(x_k-\bar x_I)$; hence (L), (D), (B), and (C) of
\cref{inh:inherited-parent} hold with the parent law
$\mathcal N(\bar x_I,1/|I|)$ and a linear, globally well-posed parent flow.
The reachable post-handoff set is four-dimensional and its centroid image
three-dimensional, so it contains distinct configurations with the same
centroid and different members, and the named-member probe has the laws
$\mathcal N(x_i,1)\neq\mathcal N(x_i',1)$ there; since every within-block
difference relaxes to a value independent of the handoff state, the
separation decays like $e^{-3t}$, and translations along the fibers are not
symmetries of the fine field, since they do not commute with that
relaxation. The full flow has events at every point of $U$ for every
$\delta\in\mathsf C$ and the nulls have none, which is (F5) in the strong
grade. The parent centroid distance relaxes at rate three to
$53/120<\delta$, so the parents unite under the descended field with the
first-stage selector, which is (F6a).

For (iii), every state of $\mathcal C$ is in the domain of the selector, so
its parent distance exceeds $\delta$, and the parent distance obeys the same
law whatever the fine state inside the blocks; it decreases to $53/120$ and
reaches $\delta$ from above.
\end{proof}

The calibration separates the clauses. With a different evaluator for each
agent the anchors supply the membership order: the full-flow order is the
centroid dendrogram of the anchors, as in \cref{inh:scalar-additive}, but the
null factor is one, so (F5) holds against both nulls, even in the strong
grade. The frozen-environment control shows that the selection already
follows from each agent's response to the initial laws of the others. What
fails is (F6b): the second stage forms from every handoff whatever the first
stage produced. Reading (F6b) with a class that contained parents within
$\delta$ would have reversed this verdict. The example concerns a declared
restricted-family flow; the full Fisher flow of \cref{eq:dyn-potential} does
not preserve the Gaussian family (\cref{prop:dyn-gaussian-noninvariance}).

The finite construction of this chapter settles inheritance for a supplied
hierarchy and makes no selection claim, and \cref{ch:boundary-closure} gives
exact inheritance through a supplied hierarchy without addressing (F2),
(F5), or (F6). Each later chapter states its own verdict by the report above;
the following summary records what the repaired clauses already decide.
\Cref{ch:fusion} gives explicit fusion events under a nonsmooth interaction.
Its fused children are identical, so (F4) fails, and whatever its (F5)
verdict for its declared selector it is at most a two-stage selection with
exact inheritance. The erasure construction of \cref{ch:smooth-erasure} has
one selection stage. Under the original smooth potential,
\cref{ch:conditional-formation} selects two parents and then their union on a
prepared common-precision family, with a partition sequence that does not
depend on the coupling, a separate control stated there. Its self costs are
nonzero, so the two nulls differ: against the label-preserving null the
agent-dependent self weights \cref{inh:null-self-row} change the membership
history on a nonempty open set already at zero coupling, while against the
renormalized null the partition sequence is reproduced. Its (F5) verdict is
therefore split and the clause fails; since the union forms from every
counterfactual handoff in the domain of its selector, (F6b) fails as well.
\Cref{ch:open-formation} gives a six-agent construction on a full open law
neighborhood with native preferential selection; its theorem establishes
neither (F5) nor (F6b), integrations of the limit field indicate that the
null flow forms the same blocks, and it is not shown to be a recursive
formation. The per-agent constructions of \cref{ch:per-agent-inheritance}
and \cref{ch:collective-local-parent} are one-stage results whose erased
phase, respectively common rotation, is gauge relative to every algebra
declared there, so their retentions fail (F4) as declared here. All the
smooth constructions supply their compression architectures. Under the
global free energy, \cref{ch:global-formation} selects first-level blocks and
then their union with a derived selector, the mode of the structure posterior
of the agents' current means, through parents that are exact and type-closed
for conjugate linear-Gaussian children; the label-preserving null and the
frozen-environment control reproduce its partition sequence, so (F5) fails,
and the union forms from every counterfactual handoff, so (F6b) fails, and
it is a two-stage selection with exact inheritance.

\openproblemheading{A recursive formation}{inh:formation-open}
The normalized per-agent construction of \cref{ch:per-agent-inheritance}
supplies a complementary one-stage result: grouping with conflicting
records under its declared selector constant, reported there over the
constants for which it holds, with a parent that preserves its declared port
dynamics, finite adaptive histories, and nonflat geometry on a finite
window. That parent quotients a hidden phase that is gauge relative to every
algebra the chapter declares and is distinguished only by a child-level law
diagnostic outside its probe rule, so it fails (F4). The phase interface is
predeclared and the parent retains primitive ports. Joining already
canonical parents loses no new coordinate, so neither a second compressive
handoff nor (F6b) follows.

The global construction of \cref{ch:global-formation} has a derived selector
with no constant beyond the constitutive data, exact and type-closed parents
for conjugate linear-Gaussian children, and a parent that acts on its
children through the expected child kernel; it fails (F5) and (F6b). With
block kernels that do not depend on the block and no multi-agent record, its
first selection event is reproduced exactly by the label-preserving null
(\cref{gf:first-event}), so interaction necessity can only come from later
stages, and with a unique self-consistent second-level structure the union
does not depend on the handoff.

Construct a population dynamics under the global free energy that
satisfies (F1) to (F6), preferably in the strong and amplifying grades, with
a derived selector. The routes left open by that chapter are block priors
that depend on the block or multi-agent records, so that the first event is
not a null event; second-level multistability reached along nested paths, so
that the union depends on the handoff; and dynamical inheritance for child
and record kernels outside the conjugate linear-Gaussian family, where the
block statistics of \cref{gf:structure-factor} still give static inheritance
but the agent update is not affine in its record (\cref{gf:beyond}). The
same problem stands for the per-agent free energies, and for the reciprocal
potential as the configuration-level device it is. Determine also whether
type closure and an action of the parent on its children should become
clauses of the standard.
